% UG54 study guide: the conceptual core, The Panel and Portfolio Mathematics through Conditional Strategies II.
% Compile with pdflatex (twice, for the references).
% To hand out the exercises without answers, change \solutionstrue to \solutionsfalse below.
\documentclass[11pt]{article}
\usepackage[margin=1in]{geometry}
\usepackage{amsmath,amssymb}
\usepackage{booktabs,array}
\usepackage{enumitem}
\usepackage{xcolor}
\usepackage[most]{tcolorbox}
\usepackage{fancyhdr}
\usepackage[hidelinks]{hyperref}

\newif\ifsolutions
\solutionstrue          % \solutionsfalse hides the Solutions section

\setlength{\parindent}{0pt}
\setlength{\parskip}{5pt}
\setlist{itemsep=2pt,topsep=3pt,leftmargin=1.4em}
\pagestyle{fancy}\fancyhf{}
\lhead{\small UG54 Data-Driven Investing}\rhead{\small Study guide: the conceptual core}
\cfoot{\small\thepage}

% the formula box
\newtcolorbox{keybox}[1][Formulas]{colback=gray!7,colframe=gray!55,boxrule=0.4pt,arc=2pt,
  left=6pt,right=6pt,top=3pt,bottom=3pt,title={#1},fonttitle=\bfseries\small,coltitle=black,colbacktitle=gray!22}

% exercises, numbered by section; \sol{label} starts the matching solution
\newcounter{ex}[section]
\renewcommand{\theex}{\thesection.\arabic{ex}}
\newenvironment{ex}[1][calculation]{\refstepcounter{ex}\par\medskip\textbf{\theex}\ \ {\small\textsc{[#1]}}\ \ }{\par}
\newcommand{\sol}[1]{\par\smallskip\textbf{\ref{#1}}\ \ }
\newcommand{\tfd}{\emph{True, False or Depends?}\ }
\newcommand{\head}[1]{\par\medskip\textbf{#1}\par}
\newcommand{\re}{r^{e}}
\newcommand{\SR}{\mathit{SR}}
\newcommand{\AR}{\mathit{AR}}
\newcommand{\IR}{\mathit{IR}}
\newcommand{\RV}{\mathit{RV}}

\begin{document}

\begin{center}
{\LARGE\bfseries UG54 Data-Driven Investing with Python}\\[4pt]
{\Large Study guide: the conceptual core}\\[4pt]
{\normalsize From \emph{The Panel and Portfolio Mathematics} to \emph{Conditional Strategies II}}\\[2pt]
{\small Prof.\ Alan Moreira \quad Fall 2026}
\end{center}

\head{What this is}
The course has two halves that run together: building things in Python with an AI, and the ideas that tell you what to
build and whether to believe the result. This guide covers the second half only: the formulas, the calculations and
the arguments you should be able to do with a pencil. It is not the whole course. The coding and the implementation are
in the notebooks.

\head{How to use it}
Each section has the same four parts: the ideas, the formulas, a worked example, and exercises. The exercises are of
three kinds. \textsc{[calculation]} asks for a number. \textsc{[concept]} asks for one or two sentences; many are
``True, False or Depends'', and as on the exams the explanation is what counts. \textsc{[code]} shows a few lines and asks
you to read or fix them, or asks you to give instructions to an AI. You will not be asked to write code.

Do the exercises before you read the solutions at the end. If the arithmetic is ugly, write down the calculation and
leave it. A formula sheet is on pages \pageref{sec:sheet}--\pageref{sec:sheetend}.

\head{Ten ideas that run through everything}
\begin{enumerate}[itemsep=1pt]
\item \textbf{Timing.} Use what you knew at the end of month $t-1$. Earn the return of month $t$. Most backtests that look
too good break this rule somewhere.
\item \textbf{A return has three parts:} $\re = \alpha + \beta f + \varepsilon$. Beta is cheap to buy, alpha is scarce, and
$\varepsilon$ is noise you should diversify.
\item \textbf{Variances add, volatilities do not:} $\sigma^{2} = \beta^{2}\sigma_{f}^{2} + \sigma_{\varepsilon}^{2}$.
\item \textbf{Squared Sharpe ratios add} when bets are uncorrelated: $\SR^{2}_{\max} = \SR_{m}^{2} + \AR^{2}$.
\item \textbf{Size a bet by mean over variance:} $x = \mu/(\gamma\sigma^{2})$. So the volatility you give a bet is proportional to its
Sharpe (or appraisal) ratio.
\item \textbf{Risk is learned fast, means slowly.} A ratio measured over $T$ years has a standard error of about
$1/\sqrt{T}$, so $t \approx \text{ratio}\times\sqrt{T}$.
\item \textbf{An in-sample maximum is algebra, not evidence.} The best of many backtests, the tangency portfolio, the published
anomaly: each was chosen because it looked good.
\item \textbf{A high expected return is a low price.} Ask who is selling to you, and why they are willing.
\item \textbf{Estimate the risk, supply the alpha, and shrink what you are unsure of.}
\item \textbf{Say exactly what you want, then check what you got.} The AI will not ask which price, which lag or which
denominator you meant.
\end{enumerate}

\tableofcontents
\newpage

% =====================================================================================================
\section{Returns, portfolios and weights}
{\small\emph{From: The Panel and Portfolio Mathematics}}

\head{The ideas}
\begin{itemize}
\item A portfolio is a list of weights. Its return is the weighted average of the returns of what it holds.
\item Weights that add to one: you invested all your money. Weights that add to zero: a long-short position that finances
itself. The excess return $r - r_{f}$ is the simplest one: $+1$ in the asset, $-1$ in T-bills.
\item Value weights \emph{are} the market: the only portfolio every investor can hold at the same time. They rebalance
themselves when prices move. Equal weights put the same dollars in a tiny firm as in a giant, so an equal-weighted
portfolio is mostly a bet on small stocks, and it must be traded back to equal every month.
\item Diversification shrinks volatility, not the tail. The average stock has a volatility of about 52\% a year; the
market, about 15.5\%. It does not go to zero because stocks move together.
\item Individual stocks have fat tails and are skewed to the right. The market is skewed to the left: crashes happen to
everything at once. The Sharpe ratio sees none of this.
\end{itemize}

\begin{keybox}
\[ r_{p,t} = \sum_{i} w_{i,t-1}\, r_{i,t}
\qquad\text{value weights: } w_{i,t-1} = \frac{ME_{i,t-1}}{\sum_{j} ME_{j,t-1}}
\qquad\text{equal weights: } w_{i} = \frac{1}{N} \]
\[ \text{compounding: } 1 + r_{1\to T} = (1+r_{1})(1+r_{2})\cdots(1+r_{T}) \]
\[ \text{monthly to yearly: mean} \times 12, \quad \text{volatility} \times \sqrt{12}, \quad \text{Sharpe ratio} \times \sqrt{12}
   \quad(\text{daily: } 252,\ \sqrt{252}) \]
Equal weights, $N$ assets, each with volatility $\sigma$, average correlation $\bar\rho$:
\[ \sigma_{p}^{2} = \frac{\sigma^{2}}{N}\big(1 + (N-1)\bar\rho\big) \;\longrightarrow\; \sigma^{2}\bar\rho \quad\text{as } N \text{ grows} \]
\end{keybox}

\head{Worked example}
Three stocks. Market values at the end of last month: \$500, \$300 and \$200 billion. Returns this month: $+4\%$, $-2\%$
and $+10\%$.

Value weights are $0.5$, $0.3$, $0.2$, so $r_{VW} = 0.5(4) + 0.3(-2) + 0.2(10) = 3.4\%$. Equal weights give
$r_{EW} = (4 - 2 + 10)/3 = 4\%$. A long-short that buys stock 3 and shorts stock 2 has weights $(0, -1, +1)$ and returns
$10 - (-2) = 12\%$ on each dollar of the long leg.

\head{Traps}
\begin{itemize}
\item Weighting by this month's market value. A stock that rose this month has a higher market value at the end of the
month, so you give more weight to the winners after the fact. In class this turned 15.7\% a year into 25.5\%.
\item Lagging a stacked panel without grouping by stock: the lag reaches into the previous company.
\item Dropping missing values \emph{after} forming the weights: the weights no longer add to one.
\item Adding returns over time instead of compounding. Multiplying a mean by $\sqrt{12}$.
\end{itemize}

\head{Exercises}

\begin{ex}\label{ex:1.1}
Four stocks have market values of \$600, \$250, \$100 and \$50 billion at the end of last month. This month they return
$+2\%$, $-4\%$, $+10\%$ and $+20\%$. (a) Compute the value-weighted and the equal-weighted return. (b) Why is the gap
so large? (c) If you hold the value-weighted portfolio and trade nothing, what are your weights at the end of the
month? Are they still value weights? What would the equal-weighted investor have to do?
\end{ex}

\begin{ex}\label{ex:1.2}
(a) A stock returns $+50\%$ and then $-50\%$. What did \$1 become? What would you have said if you had added the two
returns? (b) A strategy has a monthly average return of 1\% and a monthly standard deviation of 5\%. T-bills pay
0.25\% a month. Give the yearly mean, the yearly volatility and the yearly Sharpe ratio.
\end{ex}

\begin{ex}\label{ex:1.3}
There are 5{,}000 stocks, each with a volatility of 52\% a year. (a) If they moved independently, what would the
volatility of the equal-weighted portfolio be? (b) The actual portfolio has a volatility of 15.5\%. What average
correlation between stocks does that imply? (c) What does (b) tell you about how much risk a 10{,}000-stock portfolio
could remove?
\end{ex}

\begin{ex}[concept]\label{ex:1.4}
\tfd ``Individual stocks are skewed to the right, so a portfolio of many stocks is skewed to the right too.''
\end{ex}

\begin{ex}[concept]\label{ex:1.5}
Give two reasons the value-weighted portfolio, and not the equal-weighted one, deserves to be called ``the market''.
Then give one reason the measured return of the equal-weighted portfolio overstates what an investor could earn.
\end{ex}

\begin{ex}[code]\label{ex:1.6}
\texttt{df} has one row per stock and month, with columns \texttt{permno}, \texttt{date}, \texttt{ret} (the return during
the month) and \texttt{me} (market value at the end of the month). An AI wrote:
\begin{verbatim}
1  df['w'] = df['me'].shift(1)
2  vw = df.groupby('date').apply(lambda g: np.average(g['ret'], weights=g['me']))
\end{verbatim}
Line 1 was meant to give last month's market value. Line 2 was meant to give the value-weighted market return.
(a) What is wrong with line 1? (b) What is wrong with line 2, and in which direction is the answer biased? (c) Neither
line gives an error. How would you catch each mistake?
\end{ex}

% =====================================================================================================
\section{Sorts and long-short portfolios}
{\small\emph{From: Sorts, Breakpoints, and Long-Short Portfolios}}

\head{The ideas}
\begin{itemize}
\item The recipe never changes: \textbf{rank} every stock within the month on last month's signal, \textbf{bucket}
them, \textbf{weight} inside each bucket, \textbf{average} to get each bucket's return, and take the \textbf{spread}:
top bucket minus bottom bucket.
\item You are betting on a property, not on names. Membership turns over as firms change.
\item Each leg's weights add to one, so the long-short's weights add to zero. A book of \$100 means \$100 long and \$100
short. There is no T-bill rate to subtract.
\item Read all the buckets. Returns that rise step by step are stronger evidence than two odd extremes. Look at each
leg separately: a spread that comes only from the short leg may be one you cannot trade.
\item More buckets give a wider spread and fewer stocks in each bucket, so more noise. The $t$-statistic rises and then
falls.
\item Most stocks are small. An equal-weighted sort is mostly a sort of small stocks. The research convention is
breakpoints from NYSE stocks and value weights inside each bucket.
\end{itemize}

\head{Worked example}
Six stocks, sorted on a signal known last month. Use the top two and the bottom two.
\begin{center}
\begin{tabular}{lcccccc}
\toprule
stock & A & B & C & D & E & F \\
\midrule
signal (last month) & 0.10 & 0.15 & 0.20 & 0.30 & 0.40 & 0.50 \\
market value last month (\$bn) & 10 & 40 & 200 & 100 & 150 & 50 \\
return this month & $-3\%$ & $-1\%$ & $0\%$ & $1\%$ & $2\%$ & $4\%$ \\
\bottomrule
\end{tabular}
\end{center}
Long leg (E, F), value-weighted: $(150\times 2 + 50\times 4)/200 = 2.5\%$. Short leg (A, B):
$(10\times(-3) + 40\times(-1))/50 = -1.4\%$. The long-short earns $2.5 - (-1.4) = 3.9\%$. Equal-weighted it is
$3 - (-2) = 5\%$: the small stocks A and F had the most extreme returns, and equal weights give them more say.
A \$100 book holds \$75 of E and \$25 of F, and is short \$20 of A and \$80 of B.

\head{Exercises}

\begin{ex}\label{ex:2.1}
Use the table in the worked example, but now sort into three groups of two and report all three value-weighted returns.
(a) Is the pattern monotone? (b) The middle group holds C and D, which are more than half of the market's value. What
does that tell you about how much of the market your long-short actually trades?
\end{ex}

\begin{ex}[concept]\label{ex:2.2}
You sort stocks into ten portfolios on a signal. Portfolios 2 to 10 all earn between 7\% and 8\% a year. Portfolio 1
earns $-6\%$. The long-short, 10 minus 1, earns 14\% a year with a $t$-statistic of 4. Name two things you would check
before calling this a strategy.
\end{ex}

\begin{ex}[concept]\label{ex:2.3}
Going from 5 to 10 to 20 to 50 buckets, the spread between the top and bottom bucket keeps growing. The $t$-statistic
rises up to 20 buckets and then falls. Explain both facts.
\end{ex}

\begin{ex}[concept]\label{ex:2.4}
\tfd ``To get the Sharpe ratio of a long-short portfolio, first subtract the T-bill rate from its return.''
\end{ex}

\begin{ex}[code]\label{ex:2.5}
You ask an AI: ``Sort stocks into deciles on book-to-market and give me the return of high minus low.'' List five
things this instruction leaves open, each of which changes the answer.
\end{ex}

% =====================================================================================================
\section{Performance evaluation with one factor}
{\small\emph{From: Performance Evaluation, and Factor Models}}

\head{The ideas}
\begin{itemize}
\item Three ratios, one formula: average active return over tracking error. What changes is the benchmark. Against cash
it is the \textbf{Sharpe ratio}. Against a mandate it is the \textbf{information ratio}. Against the best copy of the
fund you can build from the factors it is the \textbf{appraisal ratio}.
\item The regression $\re_{t} = \alpha + \beta\,\re_{m,t} + \varepsilon_{t}$ is always valid; it is only statistics. The
content is in the reading: $\beta\,\re_{m}$ is what you could have bought for almost nothing, $\alpha$ is what you could not,
and $\varepsilon$ is risk nobody pays you for.
\item The same equation does two jobs. As a \emph{risk model} it uses $\beta$ and $\sigma_{\varepsilon}$. As an
\emph{expected-return model} it uses $\alpha$. A high $R^{2}$ is good news for the first job and leaves less room for the
second.
\item A significant $\alpha$ means one of two things: skill or mispricing, or a factor missing from your model.
\item Hedging removes the market: hold the fund and short $\beta$ dollars of the market for each dollar of the fund. What
is left has mean $\alpha$ and volatility $\sigma_{\varepsilon}$. It frees risk budget for the part you have a view on. It does
not always raise the Sharpe ratio, because it also removes $\beta E[\re_{m}]$ from the return.
\item If you already hold the market, what a fund adds depends on its appraisal ratio only. Its own Sharpe ratio does not
enter.
\end{itemize}

\begin{keybox}
\[ \SR = \frac{E[\re]}{\sigma(\re)} \qquad \IR = \frac{E[r_{p} - r_{b}]}{\sigma(r_{p} - r_{b})} \qquad
   \AR = \frac{\alpha}{\sigma_{\varepsilon}} \]
\[ E[\re] = \alpha + \beta\,E[\re_{m}] \qquad
   \sigma^{2} = \underbrace{\beta^{2}\sigma_{m}^{2}}_{\text{systematic}} + \underbrace{\sigma_{\varepsilon}^{2}}_{\text{its own}}
   \qquad R^{2} = \frac{\beta^{2}\sigma_{m}^{2}}{\sigma^{2}} \]
\[ \text{hedged: } \re - \beta\,\re_{m} = \alpha + \varepsilon
   \qquad \SR^{2}_{\max} = \SR_{m}^{2} + \AR^{2}
   \qquad t(\alpha) \approx \AR \times \sqrt{\text{years}} \]
\[ \text{risk budget: dollars you can hold} = \frac{\text{budget in dollars of volatility}}{\text{volatility per dollar}} \]
\end{keybox}

\head{Worked example}
A fund has an average excess return of 12\%, a volatility of 24\% and a market beta of 1.2. The market has an average
excess return of 8\% and a volatility of 16\%.
\begin{itemize}
\item $\alpha = 12 - 1.2 \times 8 = 2.4\%$.
\item Systematic volatility: $1.2 \times 16 = 19.2\%$. Its own: $\sigma_{\varepsilon} = \sqrt{24^{2} - 19.2^{2}} = 14.4\%$.
$R^{2} = 19.2^{2}/24^{2} = 0.64$.
\item $\SR = 12/24 = 0.50$, the same as the market's $8/16$. $\AR = 2.4/14.4 = 0.17$.
\item Added to the market in the best proportion: $\sqrt{0.50^{2} + 0.17^{2}} = 0.53$.
\item After 10 years, $t(\alpha) \approx 0.17\times\sqrt{10} = 0.5$. To reach $t = 2$ you would need $(2/0.17)^{2} = 144$ years.
\item A risk budget of \$6 million of yearly volatility buys $6/0.24 = \$25$ million of the fund unhedged. Hedged, it buys
$6/0.144 = \$41.7$ million, with $1.2\times 41.7 = \$50$ million of the market sold short. Expected dollars: $25\times 12\% =
\$3.0$ million unhedged, $41.7\times 2.4\% = \$1.0$ million hedged.
\end{itemize}
The hedged book earns less and has a lower Sharpe ratio (0.17 against 0.50). It can still be the right book for an
investor who already owns the market elsewhere: for them the 0.50 is nothing new, and the 0.17 is.

\head{Traps}
\begin{itemize}
\item Total returns on the left and excess returns on the right: $\alpha$ swallows the T-bill rate.
\item No intercept in the regression: $\alpha$ is forced to zero and $\beta$ is biased.
\item Annualizing a monthly $\alpha$ with $\sqrt{12}$. It is a mean: multiply by 12.
\item Reading $R^{2}$ as a grade. It says how much of the variance the factor explains, nothing about skill.
\end{itemize}

\head{Exercises}

\begin{ex}\label{ex:3.1}
A strategy has an average excess return of 9\%, a volatility of 15\% and a market beta of 0.5. The market has an average
excess return of 6\% and a volatility of 18\%. Compute (a) $\alpha$, (b) $\sigma_{\varepsilon}$ and $R^{2}$, (c) the Sharpe ratio
and the appraisal ratio, (d) the best Sharpe ratio from mixing the strategy with the market, (e) the $t$-statistic of
$\alpha$ after 9 years, and how many years it takes to reach $t = 2$.
\end{ex}

\begin{ex}\label{ex:3.2}
The market has an expected excess return of 8\% and a volatility of 16\%. Fund X has a beta of 1.2, a volatility of 20\%
and an expected excess return of 11\%. Fund Y has a beta of $-0.375$, a volatility of 10\% and an expected excess return
of 1\%. (a) Compute each fund's Sharpe ratio. (b) Compute each fund's alpha, its own volatility $\sigma_{\varepsilon}$ and
its appraisal ratio. (c) An investor holds the market and may add one fund. Which one, and what Sharpe ratio do they
reach? (d) Explain in words how a fund with a Sharpe ratio of 0.1 can be the better addition.
\end{ex}

\begin{ex}\label{ex:3.3}
Take the strategy of exercise \ref{ex:3.1}. Your risk budget is \$3 million of yearly volatility. (a) How many dollars of
the strategy can you hold unhedged? (b) Hedged, what do you trade per dollar of the strategy, and how many dollars can
you hold? (c) Expected dollars a year in each case. (d) In this example the hedged book earns fewer dollars. Give the
condition on $\alpha$, $\beta$ and the market premium under which hedging lowers expected dollars for a given risk budget.
You may leave it as an inequality.
\end{ex}

\begin{ex}[concept]\label{ex:3.4}
For each, \tfd
(a) ``A fund with an $R^{2}$ of 0.20 against the market is poorly described by the model, so its alpha cannot be trusted.''
(b) ``A value manager with an information ratio of 0.6 against the S\&P 500 has shown skill.''
(c) ``If two funds have the same alpha, an investor who holds the market should be indifferent between them.''
\end{ex}

\begin{ex}[concept]\label{ex:3.5}
A regression gives $\beta = 0.59$ and $\alpha = 6.5\%$ a year for Berkshire Hathaway, with a $t$-statistic of 1.8 over 23
years. (a) What appraisal ratio does that $t$ imply? (b) Berkshire was \emph{more} volatile than the market with a beta
well below one. Where did its risk come from? (c) Give two reasons, beyond the $t$-statistic, to be careful before
concluding that this is skill.
\end{ex}

\begin{ex}[code]\label{ex:3.6}
\texttt{r} is a fund's monthly return, \texttt{mkt} the market's monthly excess return and \texttt{rf} the monthly T-bill
rate. An AI wrote:
\begin{verbatim}
1  fit   = sm.OLS(r, mkt).fit()
2  alpha = fit.params['const'] * np.sqrt(12)
\end{verbatim}
Find three mistakes. For each, say what it does to the alpha you report.
\end{ex}

% =====================================================================================================
\section{Where factors come from}
{\small\emph{From: Where Factors Come From, and the Factor Zoo}}

\head{The ideas}
\begin{itemize}
\item Three kinds of factor model, by what you have to supply. \textbf{Time-series}: you supply the factor returns and
estimate the betas (Fama--French). \textbf{Characteristic}: you supply the exposures, the characteristics, and
estimate what each one paid (BARRA, Fama--MacBeth). \textbf{Statistical}: you supply nothing (PCA).
\item A sort makes a factor. A regression uses one.
\item Hundreds of published signals fall into a few families: value, momentum, profitability, investment, low risk,
issuance, liquidity, lead-lag. Each has two stories. \textbf{Risk}: the premium is pay for something unpleasant, so it
should last. \textbf{Mispricing}: investors make a mistake, so it should shrink once it is known and traded.
\item A price ratio has to forecast something: a high return or low growth. That is what a price is.
\item Momentum is built from returns alone: the return from month $t-12$ to month $t-2$. The signal is lagged twice, for
two reasons. The first lag makes it tradable. The second skips the month in which stocks tend to reverse.
\item A spread in average returns is not yet a new factor. It adds something only if its alpha against what you already
hold is not zero.
\item Count bets, not papers. Correlate the \emph{strategies}, not the signals. Inside a family strategies are highly
correlated; three volatility papers are one trade.
\end{itemize}

\begin{keybox}
\[ \text{Gordon: } P = \frac{D}{r - g} \quad\Longrightarrow\quad \frac{D}{P} = r - g \]
\[ \text{momentum signal for the portfolio held in month } t:\quad
   (1+r_{t-12})(1+r_{t-11})\cdots(1+r_{t-2}) - 1 \]
\[ \alpha_{\text{long-short}} = E[r_{LS}] - \beta\,E[\re_{m}]
   \qquad \SR_{\max}^{2} = \SR_{m}^{2} + \left(\frac{\alpha}{\sigma_{\varepsilon}}\right)^{2} \]
\end{keybox}

\head{Worked example}
High-signal stocks have higher betas, so a long-short has $\beta = 0.5$. The market premium is 8\%. The long-short earns
4\% a year. Its alpha is $4 - 0.5\times 8 = 0$. The signal does sort stocks by expected return, but only through market
exposure. Half a unit of the market earns the same 4\% with less risk.

It cuts both ways. A value long-short that earns 5.2\% with $\beta = -0.20$ has an alpha of
$5.2 + 0.20\times 8 \approx 6.8\%$, larger than its average return.\footnote{In class the number was 7.1\%, because the
market premium over 1980--2000 was above 8\%.}

\head{Exercises}

\begin{ex}\label{ex:4.1}
Stock A has a dividend yield of 5\%. Stock B has a dividend yield of 1\%. (a) If investors require a return of 8\% on
both, what dividend growth is priced into each? (b) If instead both are expected to grow at 3\% forever, what return
does each offer? (c) Which of the two readings is the ``value earns more'' story, and which is the ``value is cheap for
a reason'' story?
\end{ex}

\begin{ex}\label{ex:4.2}
You form a momentum portfolio on the last trading day of June 2025 and hold it through July 2025. (a) Which months of
returns go into the signal? (b) Which month is skipped, and why? (c) What happens to the backtest if you include July
2025 in the signal? (d) What would you expect if you sorted on June 2025 alone?
\end{ex}

\begin{ex}\label{ex:4.3}
The market premium is 8\%. Strategy P earns 6\% a year with a market beta of 0.75. Strategy Q earns 5\% a year with a
market beta of $-0.25$. (a) Compute each alpha. (b) An investor holds the market. What does each strategy add? (c) Which
strategy would look better in a table of average returns?
\end{ex}

\begin{ex}[concept]\label{ex:4.4}
A paper documents that a signal earned 8\% a year. After publication more investors trade it. Say what you expect to
happen to (a) the \emph{realized} return of the strategy while the money is flowing in and (b) its \emph{expected}
return afterwards, under the mispricing story. (c) How does your answer to (b) change if the premium is pay for risk?
\end{ex}

\begin{ex}[concept]\label{ex:4.5}
Two signals, illiquidity and size, have a correlation of 0.07 across stocks. The two long-short strategies built from
them have a correlation of 0.89. (a) How is that possible? (b) Which number matters for a portfolio that holds both?
(c) A report says the average correlation among 30 strategies is 0.05, so they are 30 independent bets. What is
missing?
\end{ex}

% =====================================================================================================
\section{Multi-factor models}
{\small\emph{From: Multi-Factor Models}}

\head{The ideas}
\begin{itemize}
\item Every factor in the model is a long-short portfolio that someone already sells cheaply. Alpha is the average
return that \emph{none} of the factors in the model explains.
\item Alpha belongs to a strategy \emph{and a model}. Momentum's alpha was 18\% a year against the market and 3.5\% once
UMD was in the model. Never report an alpha without naming the model.
\item Adding a factor does not always lower alpha. A negative loading on a factor that paid raises it.
\item Fix the model before you look, and justify it by what an investor can buy cheaply. Running four models and
reporting the best is not analysis.
\item Read the loadings before the alpha: a loading times the factor's premium is the part of the return that exposure
accounts for. ARKK earned 26.7\% a year over T-bills, and 21.4 points of it were a beta of 1.49 on a market that paid
14.4\%.
\item A loading says the fund \emph{co-moves} with the factor. It does not say what the fund holds.
\item Daily data sharpens the betas. It does almost nothing for the $t$-statistic of alpha, which depends on the span in
years.
\item There are two ways to estimate the same thing. \textbf{Time-series}: one regression per strategy, through time.
You supply factor returns and get betas and alpha. \textbf{Cross-sectional} (Fama--MacBeth): one regression per month,
across stocks. You supply characteristics and get what each characteristic paid that month.
\item A Fama--MacBeth slope \emph{is} a portfolio return: one unit of its own characteristic, zero of the others, zero
cost. A sort cannot do that. The intercept is not alpha; it is the return of the average stock.
\item A \textbf{characteristic-adjusted return} subtracts what the portfolio's characteristics earned. It hedges without
estimating a single beta. It needs holdings, not a track record. Its weakness: it ignores covariances.
\end{itemize}

\begin{keybox}
\[ \re_{t} = \alpha + \sum_{k}\beta_{k} f_{k,t} + \varepsilon_{t}
   \qquad E[\re] = \alpha + \sum_{k}\beta_{k}\,E[f_{k}]
   \qquad \SR_{\max} = \sqrt{\SR_{F}^{2} + \AR^{2}} \]
\[ \text{other factors uncorrelated with the market:}\qquad
   \alpha_{\text{CAPM}} = \alpha_{\text{multi}} + \sum_{k \neq m}\beta_{k}\,E[f_{k}] \]
\[ \text{hedge: per dollar of the fund, sell } \beta_{k} \text{ dollars of factor } k
   \qquad \underbrace{\beta' f_{t}}_{\text{tracking piece}} + \underbrace{\alpha + \varepsilon_{t}}_{\text{hedged piece}} \]
\[ \text{Fama--MacBeth, each month: } r_{i,t+1} = Z_{0,t} + \sum_{k} Z_{k,t}\,x_{k,i,t} + e_{i,t+1} \]
\[ \text{premium}_{k} = \text{average of } Z_{k,t} \text{ over the months}
   \qquad t = \frac{\text{mean}(Z_{k})}{\text{sd}(Z_{k})/\sqrt{\text{months}}} \]
\[ \text{one characteristic: } Z_{t} = \sum_{i} w_{i}\,r_{i}, \quad
   w_{i} = \frac{x_{i} - \bar x}{\sum_{j}(x_{j} - \bar x)^{2}}, \qquad \sum_{i} w_{i} = 0, \quad \sum_{i} w_{i}x_{i} = 1 \]
\[ \text{characteristic-adjusted return: } r_{p,t} - X_{p}'Z_{t}, \qquad X_{p} = \sum_{i} w_{i}\,x_{i} \]
\end{keybox}

\head{Worked example 1: the same fund, two alphas}
A fund earns 9.5\% a year over T-bills. Its loadings are 1.1 on the market, 0.5 on HML and $-0.4$ on UMD. The premiums
are 8\%, 4\% and 7\%.

The factors account for $1.1\times 8 + 0.5\times 4 - 0.4\times 7 = 8.8 + 2.0 - 2.8 = 8.0\%$. So the three-factor alpha is
$9.5 - 8.0 = 1.5\%$. If HML and UMD are uncorrelated with the market, the CAPM alpha is $9.5 - 8.8 = 0.7\%$. The alpha
\emph{rose} when factors were added: the fund is short momentum, a factor that paid, and still earned 9.5\%.

\head{Worked example 2: a slope is a portfolio}
One month, four stocks. The characteristic is $x = (2, 1, -1, -2)$, already centered at zero. Returns are
$(5\%, 1\%, 0\%, -2\%)$.

$\sum x^{2} = 10$, so the weights are $w = x/10 = (0.2, 0.1, -0.1, -0.2)$. The slope is
$0.2(5) + 0.1(1) - 0.1(0) - 0.2(-2) = 1.5\%$. Check: the weights add to zero, and
$\sum w_{i}x_{i} = 0.4 + 0.1 + 0.1 + 0.4 = 1$. The intercept is the equal-weighted return, 1\%.

\head{Exercises}

\begin{ex}\label{ex:5.1}
A fund earns 10\% a year over T-bills. Against a four-factor model its loadings are 0.9 on the market, 0.5 on SMB,
$-0.5$ on HML and 0.5 on RMW. The factor premiums are 8\%, 2\%, 4\% and 3\%. Assume SMB, HML and RMW are uncorrelated
with the market. (a) How much of the 10\% does each factor account for? (b) What is the four-factor alpha? (c) What
is the CAPM alpha? (d) For each dollar in the fund, what do you trade to hedge all four factors, and what do you expect
to earn on the hedged position?
\end{ex}

\begin{ex}\label{ex:5.2}
One month, four stocks. The characteristic is $x = (3, 1, -1, -3)$ and the returns are $(4\%, 2\%, -1\%, -1\%)$.
(a) Compute the weights of the slope portfolio. (b) Compute the slope. (c) Verify the two properties of the weights.
(d) What is the intercept, and which portfolio is it?
\end{ex}

\begin{ex}\label{ex:5.3}
A Fama--MacBeth regression on 240 months and about 4{,}000 stocks a month gives slopes on book-to-market with an average
of 0.4\% a month and a standard deviation of 2\% a month. (a) Compute the $t$-statistic. (b) Where do the 4{,}000 stocks
enter your calculation?
\end{ex}

\begin{ex}\label{ex:5.4}
A portfolio holds two stocks, half in each. Their size scores are $2$ and $1$; their value scores are $-1$ and $-2$
(in standard deviations from the average stock). This month one standard deviation of size paid $-0.5\%$ and one of
value paid $+1.0\%$. The portfolio returned 1\% and the equal-weighted market 2\%. (a) Compute the portfolio's size and
value scores. (b) What return did its characteristics imply, relative to the average stock? (c) What is its
characteristic-adjusted return? (d) What did you need to know about the portfolio to do this, and what did you not
need?
\end{ex}

\begin{ex}[concept]\label{ex:5.5}
For each, \tfd
(a) ``Our strategy has an alpha of 6\% a year, so it beats the market.''
(b) ``A model with more factors is a tougher test, so alpha can only fall as factors are added.''
(c) ``A loading of 1.2 on HML means the fund holds value stocks.''
(d) ``With daily instead of monthly data over the same five years, the $t$-statistic on alpha will be about
$\sqrt{21}$ times larger.''
\end{ex}

\begin{ex}[concept]\label{ex:5.6}
Fill in the table, then say in one sentence why a decile sort on book-to-market is not a pure play on value.
\begin{center}
\begin{tabular}{lll}
\toprule
 & time-series & cross-sectional \\
\midrule
one regression per & & \\
you must supply & & \\
you estimate & & \\
\bottomrule
\end{tabular}
\end{center}
\end{ex}

\begin{ex}[concept]\label{ex:5.7}
A fund is long-only, with a beta of 1.5, a loading of $-1$ on HML and an alpha of 5\% against a six-factor model. You must
sell it to an investor who judges funds by their CAPM alpha, in a period when value did well. (a) What problem do you
face? (b) Describe a package that holds this fund \emph{and other things} and would look better to that investor.
\end{ex}

% =====================================================================================================
\section{Decomposing a portfolio}
{\small\emph{From: Decomposing a Portfolio}}

\head{The ideas}
\begin{itemize}
\item A portfolio's exposure to a factor is the weighted average of the exposures of what it holds.
\item Its variance splits in two: factor risk, which you share with everyone, and specific risk, which belongs to the
names you own. Specific risk diversifies only if the weights are spread: the formula has $w_{i}^{2}$ in it. Berkshire's
Coca-Cola position, 41\% of the book, supplied 67\% of its specific variance.
\item Three ways to measure exposure, for three questions:
\begin{center}
{\small\begin{tabular}{@{}llll@{}}
\toprule
 & top-down & bottom-up & characteristics \\
\midrule
you need & the fund's returns & holdings, a beta per name & holdings, characteristics \\
tells you about & the past, averaged & today's portfolio & today's portfolio \\
reacts to a trade & eventually & at once & at once \\
main weakness & the portfolio changed & only what is disclosed & ignores covariances \\
\bottomrule
\end{tabular}}
\end{center}
\item A time-series regression measures an average over a portfolio that changed. It is today's exposure only if the
positions and their betas were stable.
\item Print the standard errors. Six factors on 60 months gives six numbers, and most of them are not different from
zero.
\item After every merge, check what the weights add to. If names drop out and you do not rescale, every score is scaled
down and nothing tells you.
\item Characteristics describe what you own, not what you paid. In 1999 the great value investor's book scored as giant
and expensive, because he bought cheap in 1988 and held.
\end{itemize}

\begin{keybox}
\[ r_{t} = \alpha + B f_{t} + \varepsilon_{t} \qquad b = B'w = \sum_{i} w_{i}\beta_{i}
   \qquad \Omega = B\,\Omega_{f}B' + \Omega_{\varepsilon} \]
\[ \underbrace{w'\Omega w}_{\text{total variance}} = \underbrace{b'\Omega_{f}\,b}_{\text{factor risk}} +
   \underbrace{\sum_{i} w_{i}^{2}\sigma_{\varepsilon,i}^{2}}_{\text{specific risk}}
   \qquad\text{one factor: } \sigma_{p}^{2} = b^{2}\sigma_{m}^{2} + \sum_{i} w_{i}^{2}\sigma_{\varepsilon,i}^{2} \]
\[ \text{equal weights, same } \sigma_{\varepsilon}:\quad \text{specific volatility} = \frac{\sigma_{\varepsilon}}{\sqrt{N}} \]
\end{keybox}

\head{Worked example}
Two stocks, weights 0.6 and 0.4. Betas 0.5 and 1.5. Specific volatilities 20\% and 30\%. Market volatility 16\%.

Portfolio beta: $b = 0.6(0.5) + 0.4(1.5) = 0.9$. Factor variance: $0.9^{2}\times 16^{2} = 207$. Specific variance:
$0.6^{2}\times 20^{2} + 0.4^{2}\times 30^{2} = 144 + 144 = 288$. Total: $495$, so the volatility is $\sqrt{495} = 22.3\%$.
Factor risk is $207/495 = 42\%$ of the variance. Each stock supplies half of the specific variance: the larger weight
on the quieter stock exactly offsets.

\head{Exercises}

\begin{ex}\label{ex:6.1}
Three stocks with weights 0.5, 0.3 and 0.2. Betas 0.8, 1.0 and 1.5. Specific volatilities 20\%, 30\% and 40\%. Market
volatility 15\%. Compute (a) the portfolio beta, (b) the factor variance and the specific variance, (c) the portfolio
volatility, (d) the share of variance that is factor risk, (e) the share of specific variance that comes from the
largest position.
\end{ex}

\begin{ex}\label{ex:6.2}
Every stock has a specific volatility of 30\%. (a) What is the specific volatility of an equal-weighted portfolio of
4, 25 and 100 stocks? (b) A 100-stock portfolio puts 40\% in one name and spreads the rest equally. Roughly, what is
its specific volatility? (c) What does this say about counting positions as a measure of diversification?
\end{ex}

\begin{ex}\label{ex:6.3}
A fund holds 16 stocks. You find characteristics for 8 of them, which carry 65\% of the book. Using the original weights
of those 8, you compute a size score of 1.64. (a) What is the size score of the matched part of the book? (b) What
single check would have caught the problem? (c) After fixing it, what are you still not measuring?
\end{ex}

\begin{ex}[concept]\label{ex:6.4}
Which of the three approaches would you use, and why? (a) You are evaluating a hedge fund that reports monthly
returns and discloses nothing else. (b) You are the risk manager of a fund that trades every day, and the CIO asks what
happens to market exposure if the largest position is sold. (c) You want to know whether a pension fund's stock
portfolio is tilted toward cheap or expensive stocks, as of today.
\end{ex}

\begin{ex}[concept]\label{ex:6.5}
A regression of Berkshire's returns on the market gives a beta of 0.78 (standard error 0.22). Adding up the betas of
its disclosed holdings gives 1.08. (a) Are the two statistically different? (b) Give two economic reasons they could
differ even with perfect estimates.
\end{ex}

% =====================================================================================================
\section{What data can teach you, and backtesting}
{\small\emph{From: Why We Backtest (slides), and Backtesting: Estimate, Tune, Test}}

\head{The ideas}
\begin{itemize}
\item Rank 49 industries on their alpha over the last three years, buy the top fifth, short the bottom fifth, hold a
year. The alphas you ranked on promised about 20\% a year. The trade delivered about 2\%. The same regressions gave
betas that carried over almost perfectly. \textbf{Beta carries over; alpha does not.}
\item The reason is how fast each thing is learned. Betas and volatilities are about how returns \emph{move}: sample
more often and months are enough. A mean, or an alpha, uses only where the price started and where it ended: only
more years help.
\item You need decades to measure an alpha that may last only years. So the data cannot be where the idea comes from.
The reason comes first; then the data tests it.
\item Everything you choose from data inflates your result. So use three samples with three jobs:
\begin{center}
\begin{tabular}{lll}
\toprule
sample & what you do with it & may you look? \\
\midrule
\textbf{Estimate} & fit: the alphas, the betas, the ranking & yes \\
\textbf{Tune} & choose among your rules & yes \\
\textbf{Test} & report & once, at the end \\
\bottomrule
\end{tabular}
\end{center}
\item Out of sample means out of \emph{time}. A random split trains on 1997 and tests on 1985.
\item The in-sample winner is the worst choice. Of 29 strategies, the single best in the estimate sample had a Sharpe
ratio of $-0.08$ in the test sample. Holding all 29 equally: 0.85.
\item Walk-forward repeats the idea every year: every month's return is earned on choices made before that month began.
\item Never compare strategies on return alone. Compare at the same risk.
\item \textbf{Leakage} is using data on a date that you could not have had that day: a signal that is not lagged; a
universe of survivors; accounting numbers dated to the quarter they describe and not to the day they were released;
prices adjusted for later splits; a universe changed after seeing a result.
\item You tried more than one thing. Twenty tries give one $|t| > 2$ for free. A higher cutoff flags less noise and
finds fewer real signals; only a longer sample improves both.
\end{itemize}

\begin{keybox}
\begin{center}
\begin{tabular}{lll}
\toprule
 & standard error & shrinks with \\
\midrule
mean & $\sigma/\sqrt{T}$ & years, $T$ \\
alpha & $\sigma_{\varepsilon}/\sqrt{T}$ & years, $T$ \\
Sharpe or appraisal ratio & $\approx 1/\sqrt{T}$; more exactly $\sqrt{(1 + \SR^{2}/2)/T}$ & years, $T$ \\
beta & $(\sigma_{\varepsilon}/\sigma_{f})/\sqrt{N}$ & observations, $N$ \\
volatility & $\approx \sigma/\sqrt{2N}$ & observations, $N$ \\
\bottomrule
\end{tabular}
\end{center}
\[ t = \text{ratio}\times\sqrt{T} \qquad\qquad
   \text{years needed: } T = \left(\frac{z}{\text{ratio} - \text{hurdle}}\right)^{2} \]
{\small $z$: 1.64 to be 95\% sure (one-sided), 1.96 for the usual two-sided 5\%. With $N$ tries (Bonferroni), test at
$5\%/N$: about 2.8 for 10 tries, 3.0 for 20, 3.1 for 29.}
\end{keybox}

\head{Worked example}
The market over 100 years: mean 8.3\%, volatility 18.4\%, Sharpe ratio 0.45.
\begin{itemize}
\item After ten years the standard error of the mean is still $18.4/\sqrt{10} = 5.8\%$.
\item To be 95\% sure the premium is above zero you need $(1.64/0.45)^{2} = 13$ years.
\item A strategy with a true ratio of 1 needs $(1.64/1)^{2} = 2.7$ years to be distinguished from zero, and
$(1.64/(1 - 0.5))^{2} = 11$ years to be distinguished from 0.5. Only the gap to the hurdle counts.
\item With three samples you need that in \emph{each} one: about 8 years in all for a ratio of 1.
\end{itemize}

\head{Exercises}

\begin{ex}\label{ex:7.1}
A strategy has a true appraisal ratio of 0.8. How many years of data do you need to be 95\% sure that its appraisal
ratio is (a) above zero, (b) above 0.5? (c) Explain why (b) is so much longer.
\end{ex}

\begin{ex}\label{ex:7.2}
A fund shows you three years of returns with a Sharpe ratio of 1.2. (a) Compute the $t$-statistic. (b) Compute the
standard error of the Sharpe ratio with both formulas in the box. (c) Is the fund's volatility needed for either
calculation? Why not? (d) What else would you ask about how this fund came to be in front of you?
\end{ex}

\begin{ex}\label{ex:7.3}
You regress an industry on the market using three years of data. The industry's own volatility is $\sigma_{\varepsilon} = 12\%$
a year and the market's is 16\%. (a) What is the standard error of alpha? (b) What is the standard error of beta with
monthly data? With daily data (756 days)? (c) Which of the two answers changed, and why?
\end{ex}

\begin{ex}\label{ex:7.4}
You tested 20 signals and the best one has a $t$-statistic of 2.6. (a) If none of the 20 were real, how many would you
expect to see with $|t| > 2$? (b) What $t$ does Bonferroni ask for? (c) The 20 signals include six versions of
volatility that correlate 0.9 with each other. Is the Bonferroni cutoff too strict or too lenient, and why?
\end{ex}

\begin{ex}[concept]\label{ex:7.5}
For each, \tfd
(a) ``I held out a random 30\% of the months as a test set and the Sharpe ratio there is close to the training Sharpe
ratio, so the strategy works out of sample.''
(b) ``A strategy with 5\% volatility has a more precisely measured Sharpe ratio than one with 20\% volatility over the
same years.''
(c) ``The strategy that did best in the estimate sample is the one most likely to do best in the test sample.''
\end{ex}

\begin{ex}[concept]\label{ex:7.6}
In the industry experiment, ranking on alphas from three years of data promised 20\% and delivered about 2\%. (a) The
alphas were estimated without bias. Why did the top fifth still disappoint? (b) Why did the betas not disappoint?
(c) Ranking on a single year of data gives a noisier alpha, yet the trade earns more than twice as much the next year.
Is alpha better measured with one year than with three? What else could be going on?
\end{ex}

\begin{ex}[code]\label{ex:7.7}
You have monthly returns for 1980--2000 and want to choose the look-back window of a momentum strategy (3, 6, 9 or 12
months) and report how well it works. Give the AI instructions: which years do what, what it should compute on each,
and what it should report. Then say what you must \emph{not} do after you see the final number.
\end{ex}

\begin{ex}[concept]\label{ex:7.8}
A backtest sorts stocks on their split-adjusted price in 1985 and finds that low-priced stocks did very well afterwards.
(a) Where is the leak? (b) State the rule that covers this and every other leak. (c) Which prices should you use for
the signal, and which for the returns?
\end{ex}

% =====================================================================================================
\section{Hiring and firing}
{\small\emph{From: Hiring and Firing}}

\head{The ideas}
\begin{itemize}
\item Putting money into a strategy and hiring a manager are the same decision. So are stopping and firing. Both are
statistical: you never see the true ratio, only an estimate with a standard error of about $1/\sqrt{T}$.
\item Three ingredients: a \textbf{goal} $a^{*}$, how \textbf{sure} you want to be, and \textbf{time}. Between the hire
line and the walk-away line you do not know yet. Time is the only lever that closes the gap, and it never closes
completely.
\item The record that got a manager hired is the estimate sample. Of 100 managers with no skill, 5 clear a 95\% bar, and
those are the ones who come to see you. \textbf{The test starts the day you hire.}
\item A rule that only looks at returns does not know how the manager got there. Buffett after 1999: the rule fires him
within about a year, just before one of his best stretches.
\item How far bad news should move you depends on what you believed before. That prior is your \emph{hit rate}: the share
of the people your process hires who turn out to be good. Same returns, different priors, opposite decisions: Citadel
and Point72 put money into Melvin in its worst week.
\item One outcome does not grade a decision.
\item A manager with no skill is not free: you gave them risk, and they added nothing to the return.
\end{itemize}

\begin{keybox}
\[ \text{go in when } \widehat{\AR}_{T} \ge a^{*} + \frac{z_{p}}{\sqrt{T}}
   \qquad\qquad \text{walk away when } \widehat{\AR}_{T} \le a^{*} - \frac{z_{p}}{\sqrt{T}} \]
{\small $z_{p}$: 0.84 for 80\%, 1.28 for 90\%, 1.64 for 95\%, 2.33 for 99\%. $T$ in years since you started the clock.}
\[ \underbrace{\frac{\pi_{T}}{1-\pi_{T}}}_{\text{odds he is good, after}} =
   \underbrace{\frac{\pi}{1-\pi}}_{\text{odds at hire}} \times \exp\!\big(T\,a\,(\widehat{\AR}_{T} - a/2)\big)
   \qquad \text{fire when } \widehat{\AR}_{T} < \frac{a}{2} - \frac{1}{a\,T}\ln\frac{\pi}{1-\pi} \]
Your book has Sharpe ratio $S$; the manager is hired to deliver $a$. Your Sharpe ratio with the manager:
\[ \text{if good: } \sqrt{S^{2} + a^{2}} \qquad\qquad \text{if no skill: } \frac{S^{2}}{\sqrt{S^{2} + a^{2}}} \]
\[ \text{hire if}\quad \pi \times \text{gain} \times \text{years the edge lasts} \;>\;
   (1-\pi) \times \text{loss} \times \text{years until you fire} \]
{\small $\pi$: the share of your hires who are good. Gain and loss: the change in your Sharpe ratio in each case. Years
until you fire an unskilled manager: about $(z_{p}/a)^{2}$.}
\end{keybox}

\head{Worked example}
Your goal is an appraisal ratio of 0.5 and you want to be 80\% sure ($z = 0.84$).
\begin{center}
\begin{tabular}{lccc}
\toprule
years of record & 1 & 4 & 16 \\
\midrule
go in if the record is above & 1.34 & 0.92 & 0.71 \\
walk away if it is below & $-0.34$ & 0.08 & 0.29 \\
\bottomrule
\end{tabular}
\end{center}
After one year the bar is nearly three times the goal, and a bad first year is not grounds for firing anyone. Four times
the record halves the gap.

Now with a prior. You hired a manager to deliver $a = 1$. Two years in, the record is zero. The record is
$\exp\big(2\times 1\times(0 - 0.5)\big) = 0.37$ times as likely from a good manager as from an unskilled one. If 80\% of
your hires are good, the odds go from 4 to $4\times 0.37 = 1.47$: a 60\% chance he is still good. If half are good, the
odds go from 1 to 0.37: a 27\% chance.

\head{Exercises}

\begin{ex}\label{ex:8.1}
Your goal is an appraisal ratio of 0.5 and you want to be 95\% sure. (a) Compute the go-in line and the walk-away line
after 4 years and after 16 years. (b) A strategy's true ratio is exactly 0.5. What will the rule tell you, most of the
time, at any horizon? (c) After 1 year a strategy shows a ratio of $-0.8$. Do you walk away?
\end{ex}

\begin{ex}\label{ex:8.2}
You hired a manager to deliver an appraisal ratio of 1. After three years the record is 0.2. (a) How much likelier is
this record from a good manager than from an unskilled one? (b) Two out of three of your hires turn out good. What are
the odds, and the probability, that this one is good? (c) Redo (b) for a process where four out of five are good.
(d) Do you fire in (b)? In (c)?
\end{ex}

\begin{ex}\label{ex:8.3}
A fund loses 53\% in a month. Suppose such a month is three times likelier from an unskilled manager than from a skilled
one. Before the month you believed the manager was skilled with probability 80\%. (a) What do you believe now? (b) An
investor with no inside view started at 50\%. What do they believe now? (c) Both investors saw the same return. Is one
of them making a mistake?
\end{ex}

\begin{ex}\label{ex:8.4}
Your book has a Sharpe ratio of 0.6. You consider a manager who claims an appraisal ratio of 0.8, uncorrelated with your
book. (a) What is your Sharpe ratio if the claim is true? If the manager has no skill? (b) How long until you would
fire an unskilled manager, at 95\%? (c) A real edge lasts five years. What share of your hires must be good for hiring
to pay? You may leave the last step as a fraction.
\end{ex}

\begin{ex}[concept]\label{ex:8.5}
\tfd ``A manager shows a ten-year record with an appraisal ratio of 0.6, a $t$-statistic of 1.9. Since this clears the
95\% one-sided bar, there is only a 5\% chance that she has no skill.''
\end{ex}

\begin{ex}[concept]\label{ex:8.6}
\emph{Capital Allocation I} says a pod shop might close a team after a 7.5\% loss. \emph{Hiring and Firing} fires on the
appraisal ratio. (a) Which of the two rules can a manager influence by taking less risk? (b) Why would the principal
want a rule in percent anyway? (c) What does each rule ignore?
\end{ex}


% =====================================================================================================
\section{Why should this work?}
{\small\emph{From: Why Should This Work?}}

\head{The ideas}
\begin{itemize}
\item Re-run 27 published signals on our sample: 6 clear the Bonferroni bar and 6 come out with the opposite sign.
There are three explanations. \textbf{Selection}: journals print what worked. \textbf{Overfitting}: the authors
searched too. \textbf{Decay}: it was real, and publication traded it away. Two say the effect never existed. One says you
are late. No amount of data separates them.
\item Count ideas, not papers. Three volatility strategies correlate 0.94 to 0.97: one trade, three citations.
\item Risk is a property of an asset \emph{for an investor}. An asset earns a premium because it pays off badly in bad
times. The CAPM is the special case where ``bad times'' means the market falling. Your bad times may be different: your
job, your house, your credit line.
\item So deviate from the market only if you differ from the average investor. Someone who works in the auto industry
should hold the market minus autos.
\item As you scale up, you become the investor who bears the risk. ``Risky for them, not for me'' is a statement about a
small position.
\item A high expected return is a low price. Somebody is selling to you at that price. The price can be low because the
business is expected to do badly (no premium there) or because the payoff is bad in states investors care about (a
premium). A backtest cannot tell you which.
\item Ask who is on the other side and what makes them willing. A \textbf{constraint} persists. An \textbf{error} gets
arbitraged away. A \textbf{preference} depends on who shows up. If you cannot name them, you are the counterparty.
\item Low volatility: investors who cannot borrow reach for return by buying risky, lottery-like stocks. That leaves
quiet stocks cheap. The counterparty's reason is a constraint. The risk is a bubble: in 1999 the strategy gave back 60\%
in twelve months. That is the story arriving, not the strategy breaking, and only a story written down beforehand lets
you say so.
\item A long-short has no capital in it, so its drawdown is undefined until you say what capital you posted. Leverage
turns a bad year into a forced exit.
\end{itemize}

\begin{keybox}
\[ \text{CAPM: } E[\re] = \beta\,E[\re_{m}] \qquad \text{so a negative beta should earn \emph{less} than T-bills} \]
A long-short at leverage $L$ (\$$L$ long and \$$L$ short per \$1 of capital):
\[ \text{capital}_{t} = \text{capital}_{t-1}\times(1 + L\,r_{t}) \]
\[ \text{drawdown} = \frac{\text{lowest capital after the peak}}{\text{capital at the peak}} - 1
   \qquad \text{wiped out as soon as } L \times \text{loss} \ge 100\% \]
\end{keybox}

\head{Worked example}
A long-short loses 20\% and then 25\%, measured per dollar of the long leg.
\begin{center}
\begin{tabular}{lcccc}
\toprule
leverage $L$ & 1 & 2 & 3 & 4 \\
\midrule
capital left & $0.80\times 0.75 = 0.60$ & $0.60\times 0.50 = 0.30$ & $0.40\times 0.25 = 0.10$ & $0.20\times 0 = 0$ \\
drawdown & $-40\%$ & $-70\%$ & $-90\%$ & wiped out \\
\bottomrule
\end{tabular}
\end{center}
One return series, four answers. Adding the two returns gives a fifth, $-45\%$. Code that returns
\texttt{(1 + r).cumprod()} has assumed $L = 1$ without saying so.

\head{Five questions before you trade a signal}
\begin{enumerate}[itemsep=1pt]
\item Who is on the other side, and what makes them willing?
\item Is their reason a constraint, a preference or an error?
\item Is the risk one \emph{you} can bear, at the size you intend to trade?
\item What would you have to see to stop believing? Name it now.
\item If the answer to 1 is ``I don't know'', you are the counterparty.
\end{enumerate}

\head{Exercises}

\begin{ex}\label{ex:9.1}
A long-short returns $+10\%$, $-30\%$, $-20\%$ and $+25\%$ in four periods. The peak is at the end of the first period.
(a) With \$1 long and \$1 short per dollar of capital, what is the drawdown at the trough? (b) Redo it at leverage 2, 3
and 4. (c) At leverage 1, how far below the peak are you after the fourth period? (d) You ask an AI for ``the worst
drawdown of this strategy''. Name three things the instruction must say.
\end{ex}

\begin{ex}[concept]\label{ex:9.2}
A published strategy had a $t$-statistic of 6 in the paper and has earned nothing since. For each of the three
explanations of the gap, say (a) whether the effect ever existed and (b) what you should do now. (c) Which one piece
of information, not in the returns, would help you choose among them?
\end{ex}

\begin{ex}[concept]\label{ex:9.3}
Take the low-volatility strategy. (a) Who is on the other side, and is their reason a constraint, a preference or an
error? (b) Does that make the premium more or less likely to last than momentum's, if momentum comes from slow reaction
to news? (c) What do you need to be able to do to harvest it? (d) What would an investor who ran it at three times
leverage have learned in 1999?
\end{ex}

\begin{ex}[concept]\label{ex:9.4}
For each, \tfd
(a) ``A stock with a volatility of 80\% is risky for every investor, so it must have a high expected return.''
(b) ``A strategy has a market beta of $-0.3$ and earns 5\% a year. The CAPM explains this: it is pay for being a hedge.''
(c) ``These stocks have high expected returns, so they are good businesses.''
(d) ``I cannot tell who is on the other side of my trade. That is fine as long as the backtest is honest.''
\end{ex}

\begin{ex}[concept]\label{ex:9.5}
Two investors look at the same signal and agree on every statistic. One trades it and the other does not, and both are
right. Give two concrete ways this can happen.
\end{ex}

\begin{ex}[concept]\label{ex:9.6}
A strategy had a Sharpe ratio of 1.4 in your estimate sample and 0.2 in the ten years after. Its worst year was a
bubble year in which it lost 60\%. (a) Give the reading under which you drop the strategy. (b) Give the reading under
which you keep it, and say what you change. (c) What would you have needed to write down, and when, to tell the two
apart?
\end{ex}

% =====================================================================================================
\section{Momentum}
{\small\emph{From: Momentum}}

\head{The ideas}
\begin{itemize}
\item Every month, rank stocks on their return from $t-12$ to $t-2$. Buy the top tenth, short the bottom tenth. UMD is
Ken French's version (top 30\% minus bottom 30\%).
\item The pattern in past returns: last month reverses, the past year continues, the past five years reverse. That is
why the signal skips month $t-1$.
\item UMD, 1927--2026: 7.4\% a year with a market beta of $-0.30$. So its alpha, 9.9\%, is \emph{larger} than its average
return. The losers are stocks that fell with the market; short them and you are short the market.
\item It is not a risk premium in the usual sense. A premium is pay for losing in bad times, and UMD earns $+1.3\%$ a
month in the worst third of market months. The likely story is slow reaction to news.
\item Since Carhart (1997) a new strategy has to beat momentum too. Regress it on the market \emph{and} UMD. A loading
near 1 with an alpha near 0 means it is momentum under another name. Six-month momentum is UMD again. Earnings
momentum is a different trade. Industry momentum is mostly the stock trade: its alpha falls from 13.2\% to 2.9\% once
UMD is in.
\item Why has it not been traded away? It crashes: $-65\%$ in 1932, $-53\%$ in 2009. After a bear market the losers are
the high-beta stocks, so momentum's beta is about $-0.7$, and it is short the rebound.
\item The crashes come when volatility is high, and volatility is predictable. Hold less after a volatile month. Sharpe
ratio 0.45 to 0.94; worst drawdown $-78\%$ to $-33\%$. No return forecast is needed.
\item The price of the managed version: leverage (above 2 in a quarter of months), trading, and a constant chosen with
hindsight.
\end{itemize}

\begin{keybox}
\[ UMD_{t} = \alpha + \beta\,\re_{m,t} + \varepsilon_{t} \qquad \alpha = E[UMD] - \beta\,E[\re_{m}] \]
\[ \text{is it more than UMD?}\qquad r_{t} = \alpha + \beta_{m}\,\re_{m,t} + \beta_{u}\,UMD_{t} + \varepsilon_{t} \]
\[ \text{volatility-managed:}\qquad UMD^{\sigma}_{t} = \frac{c}{\RV_{t-1}}\,UMD_{t}
   \qquad \RV_{t-1} = \sum_{d\,\in\,t-1}(r_{d} - \bar r)^{2} \]
{\small $c$: one constant, chosen so the managed series is as volatile as UMD. It sets the size, not the Sharpe ratio.}
\end{keybox}

\head{Worked example}
UMD earns 7.4\% a year with $\beta = -0.30$. The market premium is 8.3\%. So $\alpha = 7.4 - (-0.30)(8.3) = 9.9\%$.

After a two-year bear market UMD's beta is $-0.68$. If the market then rebounds 20\% in a month, the beta alone costs
$0.68\times 20 = 13.6\%$ in that month.

If momentum's daily volatility doubles from one month to the next, its variance is four times higher, and the managed
position for the following month is a quarter of what it was.

\head{Exercises}

\begin{ex}\label{ex:10.1}
The market premium is 8\% and UMD's is 7\%. (a) A ``new'' strategy earns 12\% a year, with loadings of 0.1 on the market
and 1.2 on UMD. What is its alpha against the two? What do you conclude? (b) Another strategy earns 10\% a year with
loadings of 0 and 0.2. Same questions. (c) Both strategies have large CAPM alphas. Why is the CAPM alpha the wrong
number for an investor who can buy a momentum fund?
\end{ex}

\begin{ex}\label{ex:10.2}
You run volatility-managed momentum with $c$ set so that the weight is 1 when momentum's daily volatility last month
was 1\%. (a) What is the weight after a month with a daily volatility of 0.5\%? After a month with 2\%? (b) By what
factor does the weight differ between those two months? (c) You double $c$. What happens to the average return, the
volatility and the Sharpe ratio of the managed strategy?
\end{ex}

\begin{ex}\label{ex:10.3}
After a two-year bear market UMD has a market beta of $-0.68$. After a bull market it is about zero. (a) The market
rises 15\% in a month after a bear market. What does UMD's beta contribute to its return that month? (b) Why is the beta
negative after a bear market? Think about which stocks are the losers. (c) Daniel and Moskowitz call momentum ``short a
call option on the market''. Explain the phrase using (a) and (b).
\end{ex}

\begin{ex}[concept]\label{ex:10.4}
(a) Why is momentum's alpha larger than its average return? (b) Give two reasons momentum does not look like pay for
bearing risk. (c) If it is a mistake, why has it not been traded away?
\end{ex}

\begin{ex}[concept]\label{ex:10.5}
For each, \tfd
(a) ``Industry momentum has a CAPM alpha of 13\% a year over a century, so it is a second source of alpha next to stock
momentum.''
(b) ``To run volatility-managed momentum you need a forecast of momentum's return next month.''
(c) ``The constant $c$ uses the whole sample, so the managed strategy's Sharpe ratio is overstated.''
\end{ex}

\begin{ex}[code]\label{ex:10.6}
An AI built a momentum signal with this line, where \texttt{cumret11} is each stock's cumulative return over the eleven
months ending in the current row:
\begin{verbatim}
df['signal'] = df.groupby('permno')['cumret11'].shift(1)
\end{verbatim}
(a) Is the resulting strategy valid, in the sense of having no look-ahead? (b) How does it differ from the standard
construction, and what do you expect that to do to the average return? (c) What would happen with \texttt{shift(0)}?
\end{ex}

% =====================================================================================================
\section{Capital allocation I: the frontier, and whose money it is}
{\small\emph{From: Capital Allocation I}}

\head{The ideas}
\begin{itemize}
\item Everyone wants the steepest line from T-bills: the tangency portfolio, the one with the highest Sharpe ratio.
\item Two decisions, and they separate. \emph{How to spread the money}: the tangency portfolio, the same for everyone.
\emph{How much risk}: where you sit on the line, different for everyone. A cautious investor does not hold safer assets.
They hold less of the same portfolio.
\item One asset: $x^{\star} = \mu/(\gamma\sigma^{2})$. More return, hold more; more variance, hold less. Multiply by
$\sigma$: the volatility you run is your Sharpe ratio over your risk aversion.
\item Many assets: the same formula with a matrix, $W^{\star} = \frac{1}{\gamma}\Sigma^{-1}\mu$. The minimum-variance
portfolio is $\Sigma^{-1}\mathbf{1}$; it never looks at $\mu$. Every portfolio on the frontier is a mix of the two.
\item A goal only sets the size. A risk aversion, a volatility target and a return target all give a multiple of the same
mix, with the same Sharpe ratio.
\item A new asset moves the frontier only through its \textbf{alpha} against what you already hold. A positive Sharpe
ratio is not enough: HML has a Sharpe ratio of 0.26 and an alpha of 0.2\% against the other five factors, so in sample
it adds nothing. If alpha is zero the asset is \emph{spanned}. This is the test you have run all semester.
\item The clean split rests on three assumptions. All your wealth is in financial markets. You choose for yourself.
You know $\mu$ and $\Sigma$. Each one fails.
\item \textbf{The principal.} Nobody knows their $\gamma$. Everybody can answer: the market fell 50\% in 2007--09; how much
could you stand to lose?
\item \textbf{Your job is a position.} Human capital is an asset you cannot sell, and it has betas. Hold the tangency
portfolio sized for your \emph{total} wealth, minus what looks like your job. A safe job is like owning bonds: same mix,
more risk in the account. The engineer whose pay rises with AI should hold \emph{less} Tech; the translator whom AI may
replace should hold more. People do the opposite.
\item \textbf{The delegate.} It is someone else's money, so the risk is set by a mandate: a volatility target, a
stop-loss, a tracking-error budget, limits on what you can hold. A stop in percent is a limit on volatility: at a Sharpe
ratio of 1, run a volatility about equal to your stop.
\item Every mandate has a price: the slope you lose. No borrowing costs more than no shorting. An exclusion costs
nothing or a lot, depending on what the optimizer wanted. Market-neutral gives up the market premium, on purpose: the
principal already owns the market.
\end{itemize}

\begin{keybox}
\[ \max_{x}\; x\mu - \frac{\gamma}{2}x^{2}\sigma^{2} \quad\Longrightarrow\quad
   x^{\star} = \frac{1}{\gamma}\frac{\mu}{\sigma^{2}} \qquad\qquad x^{\star}\sigma = \frac{\SR}{\gamma} \]
\[ W^{\star} = \frac{1}{\gamma}\Sigma^{-1}\mu \qquad \text{tangency} \propto \Sigma^{-1}\mu \qquad
   \text{minimum variance} \propto \Sigma^{-1}\mathbf{1} \qquad \SR_{T} = \sqrt{\mu'\Sigma^{-1}\mu} \]
\[ \text{uncorrelated assets:}\qquad w_{i} \propto \frac{\mu_{i}}{\sigma_{i}^{2}} \qquad
   \SR_{\text{tangency}} = \sqrt{\SR_{1}^{2} + \SR_{2}^{2} + \cdots} \]
\[ x \text{ times the tangency portfolio } T:\quad \text{mean} = x\,\mu_{T}, \quad
   \text{volatility} = x\,\sigma_{T}, \quad \text{same Sharpe ratio} \]
\[ \text{the alpha test:}\quad R_{\text{new}} = \alpha + \beta R_{T} + \varepsilon \qquad
   \SR^{2}_{\text{with}} = \SR^{2}_{\text{without}} + \left(\frac{\alpha}{\sigma_{\varepsilon}}\right)^{2} \]
\[ \text{with a job:}\quad w = (1+h)\,\frac{1}{\gamma}\Sigma^{-1}\mu \;-\; h\,b \qquad h = \frac{\text{human capital}}{\text{financial wealth}} \]
\[ \text{the price of a mandate, a year, at volatility } \sigma:\quad (\SR_{\text{free}} - \SR_{\text{mandate}})\times\sigma \]
{\small $w$: weights in the account. $b$: the job's betas on the assets (slopes from regressing its return on theirs).}
\end{keybox}

\head{Worked example 1: two uncorrelated assets}
Asset A has $\mu = 8\%$ and $\sigma = 20\%$. Asset B has $\mu = 3\%$ and $\sigma = 10\%$. They are uncorrelated.
\begin{itemize}
\item $\mu/\sigma^{2}$ is $0.08/0.04 = 2$ for A and $0.03/0.01 = 3$ for B. Scaled to add to one: 40\% in A, 60\% in B.
\item Mean: $0.4(8) + 0.6(3) = 5\%$. Variance: $0.4^{2}(20^{2}) + 0.6^{2}(10^{2}) = 64 + 36 = 100$, so the volatility is 10\%.
\item Sharpe ratio: $5/10 = 0.5$. The shortcut gives the same: $\sqrt{0.4^{2} + 0.3^{2}} = 0.5$.
\item B has the lower Sharpe ratio (0.3 against 0.4) and the larger weight, because it is quieter. In volatility the split
follows the Sharpe ratios: A gets $0.4\times 20 = 8$ and B gets $0.6\times 10 = 6$, a ratio of 4 to 3.
\item Sizes. With $\gamma = 2.5$: $W = (2, 3)/2.5 = (0.8, 1.2)$, twice the tangency portfolio, volatility 20\%, which is
$\SR/\gamma = 0.5/2.5$. A 15\% volatility target is 1.5 times the tangency portfolio. A return target of 8\% is 1.6 times,
with a volatility of 16\%. The Sharpe ratio is 0.5 every time.
\end{itemize}

\head{Worked example 2: a job}
Ignoring the job, the best portfolio puts 19.5\% of total wealth in Tech. Human capital equals financial wealth, so
$h = 1$. A software engineer's pay moves like $+0.5$ of Tech; a translator's like $-0.5$.
\begin{itemize}
\item Safe job ($b = 0$): the account holds $(1+1)\times 19.5 = 39\%$ in Tech.
\item Engineer: $39 - 1\times 50 = -11\%$. The job already holds more Tech than she wants.
\item Translator: $39 + 50 = 89\%$. The account is insurance: if AI wins, the job loses and the account pays.
\end{itemize}

\head{Exercises}

\begin{ex}\label{ex:11.1}
The market has an expected excess return of 6\% and a volatility of 15\%. (a) How much of your wealth do you put in the
market if $\gamma = 2$? If $\gamma = 4$? (b) What volatility do you run in each case? Check it against $\SR/\gamma$.
(c) A client cannot tell you a $\gamma$, but says that if the market fell 50\% they could stand to lose 20\% of their
wealth. How much should they hold, what volatility is that, and what $\gamma$ does it imply?
\end{ex}

\begin{ex}\label{ex:11.2}
Asset A has $\mu = 7.2\%$ and $\sigma = 12\%$. Asset B has $\mu = 6.4\%$ and $\sigma = 8\%$. They are uncorrelated.
(a) Compute the tangency weights, scaled to add to one. (b) Compute the tangency portfolio's mean, volatility and
Sharpe ratio, and check the Sharpe ratio with the shortcut. (c) A client wants 10\% volatility. What do they hold in A,
in B and in T-bills, and what return do they expect? (d) Another client wants an expected return of 8\%. What
volatility must they accept? (e) Compare the Sharpe ratios in (c) and (d), and say in one sentence why. (f) What does a
client with $\gamma = 5$ hold?
\end{ex}

\begin{ex}\label{ex:11.3}
The best portfolio of the assets you hold has a Sharpe ratio of 0.9. (a) A new asset has an alpha of 3\% against that
portfolio and a residual volatility of 10\%. What is the best Sharpe ratio once you add it? (b) Another new asset has a
Sharpe ratio of 0.5 on its own and an alpha of zero. Same question. (c) HML's alpha against the other five factors is
$0.2\%$ a year with a standard error of 1.2\%. A colleague concludes that value is dead and should be dropped. Give two
reasons to be careful.
\end{ex}

\begin{ex}\label{ex:11.4}
Ignoring your job, you would put 30\% of your total wealth in the market. Your human capital is three times your
financial wealth. How much of your \emph{account} goes into the market if (a) you are a tenured professor, with a job
that does not move with the market; (b) you are a banker, and each dollar of your human capital moves like 20 cents of
the market; (c) you have just retired? (d) What practical limit does (a) run into, and for whom does it bind the most?
\end{ex}

\begin{ex}\label{ex:11.5}
Without restrictions a client's best portfolio has a Sharpe ratio of 0.80 and a volatility of 10\%. (a) An exclusion
list lowers the best Sharpe ratio to 0.65. The client runs 12\% volatility. What does the list cost in expected return a
year? (b) The client wants 18\% volatility. With borrowing, what do they hold and what do they expect to earn? (c) The
mandate forbids borrowing. What must they do to reach 18\% volatility, and why is the result worse?
\end{ex}

\begin{ex}[concept]\label{ex:11.6}
For each, \tfd
(a) ``A cautious investor should hold safer assets than an aggressive one.''
(b) ``If every expected return doubles, the minimum-variance portfolio changes.''
(c) ``A software engineer understands technology firms, so she should hold more of them than the average investor.''
(d) ``An account whose Sharpe ratio is below that of the tangency portfolio is a mistake.''
\end{ex}

\begin{ex}[concept]\label{ex:11.7}
(a) You hand the formula $\gamma = 1$ and it tells you to hold 3.7 times your wealth in the market. Is the formula
wrong? (b) A pod shop halves a team's capital after a 5\% loss. Why does the principal impose a rule in percent
instead of telling the manager their $\gamma$? (c) Two managers have the same 5\% stop. One has a Sharpe ratio of 1,
the other of 2. Which can run more volatility, and why?
\end{ex}

% =====================================================================================================
\section{Capital allocation II: factor bets and alpha bets}
{\small\emph{From: Capital Allocation II}}

\head{The ideas}
\begin{itemize}
\item A \textbf{factor bet} is exposure anyone can buy: the market, value, Tech. It interacts with your job, your house
and your mandate, so you follow the formula loosely and tilt. An \textbf{alpha bet} is what is left after you hedge the
factors. If it is truly alpha it interacts with nothing you own, so you maximize the Sharpe ratio. The catch is the
word \emph{true}.
\item The two are estimated differently. A factor's mean is pay for risk: slow-moving, use all the history you have,
and the prior is positive. An alpha is mispricing or an edge: it decays, the prior is zero, and the evidence is often
not returns at all. Across 49 industries, alpha over five years has a correlation of $-0.12$ with alpha over the next
five. For beta it is $+0.61$; for volatility $+0.75$.
\item Hedge the factor out of each strategy and treat the hedged pieces as the assets: their means are the alphas and
their covariance is $\Sigma_{\varepsilon}$. It is the same formula.
\item If the hedged pieces are uncorrelated, the weight is $\alpha_{i}/\sigma_{\varepsilon,i}^{2}$, so \textbf{the volatility
you give a strategy is proportional to its appraisal ratio}. A bigger alpha does not mean a bigger position.
\item Correlation is an angle. Uncorrelated bets are at right angles, and squared Sharpe ratios add (Pythagoras).
Positive correlation adds less; negative correlation adds more. Judge a strategy by what it adds to what you hold.
\item Breadth: $N$ uncorrelated pods, each with a Sharpe ratio too low to be a fund, make a very good fund. Correlation
sets the limit. Our 29 strategies: 0.40 apart, 1.38 together, not $\sqrt{29}\times 0.40 = 2.13$. The scarce thing is an
\emph{uncorrelated} strategy.
\item The formula needs two inputs and you treat them differently. \textbf{Estimate the risk. Supply the alpha.} The
alpha comes from a track record, a signal or a view.
\item A hunch is not an input until it is a number: compared to what, by how much, by when, and how sure. Leaving it out
is also a number: zero. Write it down before the event, so it can be scored.
\item Your alpha is an estimate, so shrink it. The less sure you are, the smaller the bet (Black--Litterman).
\end{itemize}

\begin{keybox}
\[ r_{i} = \alpha_{i} + \beta_{i}f + \varepsilon_{i} \qquad \text{hedged: } h_{i} = r_{i} - \beta_{i}f, \quad
   \text{mean } \alpha_{i}, \quad \text{volatility } \sigma_{\varepsilon,i} \]
\[ W^{\star} \propto \Sigma_{\varepsilon}^{-1}\alpha \qquad \text{uncorrelated: } w_{i} \propto \frac{\alpha_{i}}{\sigma_{\varepsilon,i}^{2}}
   \ \Longrightarrow\ w_{i}\,\sigma_{\varepsilon,i} \propto \AR_{i}
   \qquad \SR_{\text{book}} = \sqrt{\textstyle\sum_{i}\AR_{i}^{2}} \]
Two bets with Sharpe ratios $s_{1}$, $s_{2}$ and correlation $\rho$:
\[ \SR^{2} = \frac{s_{1}^{2} + s_{2}^{2} - 2\rho\,s_{1}s_{2}}{1 - \rho^{2}} \qquad\qquad
   \text{if } s_{1} = s_{2} = s:\quad \SR = s\sqrt{\frac{2}{1+\rho}} \]
$N$ pods with average Sharpe ratio $\bar s$ and average correlation $\bar\rho$:
\[ \SR_{\text{pool}} = \bar s\sqrt{\frac{N}{1 + (N-1)\bar\rho}} \;\longrightarrow\; \frac{\bar s}{\sqrt{\bar\rho}} \text{ as } N \text{ grows}
   \qquad\qquad \bar\rho = 0:\quad \SR_{\text{pool}} = \sqrt{N}\,\bar s \]
\[ \text{Fundamental Law:}\quad \IR = IC \times \sqrt{\text{breadth}} \]
A view, as a number:
\[ \alpha = (p - (1-p)) \times \text{move} \qquad IC = \frac{\alpha}{\sigma} \text{ over the horizon}
   \qquad \AR \text{ a year} = IC\times\sqrt{\text{horizons in a year}} \]
\[ \text{volatility given} = \frac{\AR}{\SR_{\text{book}}}\times\text{book volatility}
   \qquad\qquad \text{shrink it:}\quad \alpha_{\text{act}} = \frac{\tau^{2}}{\tau^{2} + \omega^{2}}\,\alpha_{\text{view}} \]
{\small $IC$: the correlation of your forecasts with what happened. Breadth: independent bets a year. $p$: the chance
you are right. $\tau$: how big alphas tend to be. $\omega$: the standard error of your view.}
\end{keybox}

\head{Worked example 1: bigger alpha, smaller position}
\begin{center}
\begin{tabular}{lccccc}
\toprule
 & $\alpha$ & $\sigma_{\varepsilon}$ & appraisal ratio & weight $\alpha/\sigma_{\varepsilon}^{2}$ & volatility given \\
\midrule
Strategy A & 4\% & 4\% & 1.0 & $0.04/0.0016 = 25$ & $25\times 4\% = 1.0$ \\
Strategy B & 8\% & 16\% & 0.5 & $0.08/0.0256 = 3.1$ & $3.1\times 16\% = 0.5$ \\
\bottomrule
\end{tabular}
\end{center}
B has twice the alpha. It gets half the risk and an eighth of the weight. Together, if uncorrelated, the book has a
Sharpe ratio of $\sqrt{1.0^{2} + 0.5^{2}} = 1.12$. The formula divides by a variance, so it loves whatever looks quiet.

\head{Worked example 2: from a hunch to a position}
You think Apple's earnings will beat the consensus. Past beats of this size moved the stock about 4\%. You expect to be
right 60\% of the time, and being wrong means a miss. The report is in 30 trading days. Apple's stock-specific volatility
is 25\% a year.
\begin{itemize}
\item $\alpha = (0.6 - 0.4)\times 4\% = 0.8\%$ over six weeks, on Apple hedged for the market and its industry.
\item Risk over the horizon: $25\%\times\sqrt{30/252} = 8.6\%$. So the IC you are claiming is $0.8/8.6 = 0.09$. Good
managers run 0.02 to 0.05.
\item Appraisal ratio a year: $0.09\times\sqrt{252/30} = 0.27$.
\item In a book that runs 10\% volatility with a Sharpe ratio of 1, this bet gets $(0.27/1)\times 10\% = 2.7\%$ of
volatility, which is a position of $2.7/25 = 11\%$.
\item At 55\% instead of 60\%, alpha halves and so does the position. Being honest about your confidence is the formula
working.
\end{itemize}

\head{Exercises}

\begin{ex}\label{ex:12.1}
Three hedged strategies have alphas of 2\%, 6\% and 3\% and residual volatilities of 4\%, 12\% and 12\%. Their residuals
are uncorrelated. (a) Compute the three appraisal ratios. (b) Which strategy gets the largest weight? Which get the most
volatility? (c) What is the Sharpe ratio of the book? (d) The book runs 10\% volatility. How much volatility does each
strategy get, and what is each position? (e) The market has a Sharpe ratio of 0.5. What is the Sharpe ratio of the book
and the market together, and why is the formula exact here?
\end{ex}

\begin{ex}\label{ex:12.2}
Two strategies each have a Sharpe ratio of 0.6. Compute the Sharpe ratio of the best combination when their correlation
is (a) 0, (b) 0.5, (c) $-0.5$. (d) Which of the three second strategies would you pay the most for, and what does that
say about judging a strategy by its own Sharpe ratio?
\end{ex}

\begin{ex}\label{ex:12.3}
A firm runs 16 pods, each hedged, each with a Sharpe ratio of 0.5. (a) What is the firm's Sharpe ratio if the pods are
uncorrelated? (b) If their average correlation is 0.1? (c) With an average correlation of 0.1, what is the most the firm
could reach with any number of pods? (d) What should the firm look for in its next hire?
\end{ex}

\begin{ex}\label{ex:12.4}
You expect a firm to beat earnings in one month (21 trading days). A beat or a miss moves the stock about 6\%. You think
you are right 55\% of the time. The stock's specific volatility is 32\% a year. (a) What is your alpha over the month?
(b) What is the risk over the month, and what IC are you claiming? (c) What appraisal ratio a year is that? (d) Your
book runs 10\% volatility with a Sharpe ratio of 1. How large is the position? (e) Why is the alpha not ``7.2\% a
year''?
\end{ex}

\begin{ex}\label{ex:12.5}
Alphas of hedged strategies tend to be about 2\% in size ($\tau = 2\%$). You hold the view that a strategy's alpha is
5\%. What alpha do you act on if the standard error of your view is (a) 4\%, (b) 2\%? (c) What happens to the alpha you
act on, and to your position, as you become very unsure? As you become certain?
\end{ex}

\begin{ex}\label{ex:12.6}
(a) A stock-picker has an IC of 0.05 and makes 400 independent bets a year. What information ratio does the
Fundamental Law give? (b) A macro trader has an IC of 0.10 and makes 4 independent calls a year. Same question.
(c) The stock-picker's 400 bets are all on technology stocks that move together. What happens to the calculation?
\end{ex}

\begin{ex}[concept]\label{ex:12.7}
For each, \tfd
(a) ``Of two hedged strategies, the one with the larger alpha should get the larger position.''
(b) ``A strategy with a Sharpe ratio of 0.3 is not worth adding to a book with a Sharpe ratio of 1.5.''
(c) ``The average correlation among our strategies' residuals is 0.03, so they are independent bets.''
(d) ``I have no view on this stock, so my analysis makes no assumption about its alpha.''
\end{ex}

\begin{ex}[concept]\label{ex:12.8}
The same engineer considers two trades: buying a Tech fund, and buying one chipmaker against another, hedged. (a) Which
is a factor bet and which an alpha bet? (b) Which one should she size by the formula, and which one less than the
formula says? Why? (c) What does she have to believe for each trade to make sense? (d) For which input would you use 40
years of data, and for which the last year only?
\end{ex}

% =====================================================================================================
\section{Capital allocation III: estimation error}
{\small\emph{From: Capital Allocation III}}

\head{The ideas}
\begin{itemize}
\item In sample the optimizer always wins. $\Sigma^{-1}\mu$ \emph{is} the maximum in the sample it was computed on. That
is algebra, not evidence. Optimization is in-sample fitting, and everything from \emph{Backtesting} applies.
\item The formula multiplies errors in the means by $\Sigma^{-1}$. With two correlated assets the amplifier is
$(1+\rho)/(1-\rho)$: 1 at $\rho = 0$, 19 at $\rho = 0.9$. A gap of one percentage point in expected returns, which
nobody can measure, turns 50/50 into 109/$-9$.
\item Why: any pair is two bets, a level (both) and a spread (one against the other). With high correlation the spread
has a tiny variance, and the optimizer sizes by mean over variance, so a tiny gap earns a huge position.
\item Three inputs, three standard errors. Means are the worst by far: HML's mean is $2.8\% \pm 1.6\%$, while its
volatility is $10.9\% \pm 0.3\%$. Daily data helps volatilities and correlations. It does nothing for a mean.
\item Which input does the damage depends on what you hold. In a small book, the means. In a large book (29 strategies
have 406 correlations) the correlations can take over. Volatilities never do.
\item Each simple rule replaces an estimate with a constant, and is exposed only to what it still estimates:
\begin{center}
{\small\begin{tabular}{@{}llll@{}}
\toprule
rule & weights & the constant & still estimates \\
\midrule
mean-variance & $\Sigma^{-1}\mu$ & none & means, vols, correlations \\
no correlations & $\mu_{i}/\sigma_{i}^{2}$ & every correlation is zero & means, vols \\
same Sharpe & $\Sigma^{-1}\sigma$ & every Sharpe ratio is equal & vols, correlations \\
same mean (min.\ variance) & $\Sigma^{-1}\mathbf{1}$ & every mean is equal & vols, correlations \\
risk parity & $1/\sigma_{i}$ & same Sharpe, no correlations & vols \\
$1/N$ & $1/N$ & all of that, and the same vol & nothing \\
\bottomrule
\end{tabular}}
\end{center}
\item None is simply better. Each is a claim about your assets. Whether a constant helps is a race between how
different the assets really are ($\SR^{2}$) and how noisy the estimates are ($N/Y$: assets over years of data).
\item So ``$1/N$ beats the optimizer'' is a result about particular assets. Out of sample, $1/N$ wins on six factors
(0.92 against 0.83), on sectors and on asset classes. On the 29 long-shorts, which really are different, the optimizer
wins by a wide margin.
\item The race in class had one more lesson: written with the month it earns inside the estimate, mean-variance
``wins'' (1.03). One index flips the conclusion.
\item Shrinkage makes the choice continuous: stand between the estimate and the constant. Similar assets: near the
constant. Very different assets: near the estimates. A constraint (long-only, a position cap) is a prior too. The dial
is a parameter: tune it without the future, test it, report once.
\end{itemize}

\begin{keybox}
\[ \Delta W = \frac{1}{\gamma}\Sigma^{-1}\Delta\mu \qquad\qquad
   \text{two assets, same volatility:}\quad w_{1} = \frac{1}{2} + \frac{\text{gap}}{4\times\text{average mean}}\times\frac{1+\rho}{1-\rho} \]
\[ \frac{\operatorname{Var}(\text{level})}{\operatorname{Var}(\text{spread})} = \frac{1+\rho}{1-\rho}
   \qquad\qquad \text{asset 2 is shorted when } \SR_{2}/\SR_{1} < \rho \]
\begin{center}
\begin{tabular}{lll}
\toprule
input & standard error & more data means \\
\midrule
expected return & $\sigma/\sqrt{T}$ & more years \\
volatility & $\sigma/\sqrt{2N}$ & more observations (daily helps) \\
correlation & $(1-\rho^{2})/\sqrt{N}$ & more observations (daily helps) \\
\bottomrule
\end{tabular}
\end{center}
\[ \underbrace{\SR^{2}_{\text{in sample}}}_{\text{what you see}} \approx \underbrace{\SR^{2}}_{\text{real}} +
   \underbrace{N/Y}_{\text{estimation error}}
   \qquad\qquad \frac{\SR_{\text{you earn}}}{\SR} \approx \frac{1}{\sqrt{1 + \dfrac{N}{Y\,\SR^{2}}}} \]
\[ \text{shrinkage:}\quad s_{i}(d) = (1-d)\,\hat s_{i} + d\,\bar s \qquad d \approx \frac{\text{noise}}{\text{dispersion you see}} \]
{\small $N$ assets, $Y$ years of data, Sharpe ratios a year. With $N$ assets there are $N$ means, $N$ volatilities and
$N(N-1)/2$ correlations.}
\end{keybox}

\head{Worked example}
Two assets with the same volatility. The average of their expected returns is 8\%, and you estimate that asset 1 earns
one percentage point more than asset 2.
\begin{center}
\begin{tabular}{lccccc}
\toprule
correlation & 0 & 0.5 & 0.8 & 0.9 & 0.95 \\
\midrule
amplifier $(1+\rho)/(1-\rho)$ & 1 & 3 & 9 & 19 & 39 \\
weight on asset 1 & 53\% & 59\% & 78\% & 109\% & 172\% \\
\bottomrule
\end{tabular}
\end{center}
At $\rho = 0.9$: $w_{1} = 0.5 + \dfrac{0.01}{4\times 0.08}\times 19 = 0.5 + 0.59 = 1.09$. Two good assets, and the formula shorts
one of them.

The race, for the six factors over 46 years: the best in-sample Sharpe ratio is 1.17 and $1/N$ has 1.03, so the extra
$\SR^{2}$ you see is $1.17^{2} - 1.03^{2} = 0.31$. Noise alone would produce about $(N-1)/Y = 5/46 = 0.11$ of it.

\head{Exercises}

\begin{ex}\label{ex:13.1}
Two assets have the same volatility. The average of their expected returns is 6\%, and you estimate that asset 1 earns
one percentage point more than asset 2. Compute the weight on asset 1 when the correlation is (a) 0, (b) 0.5, (c) 0.8,
(d) 0.9. (e) At which correlation is asset 2 first sold short? (f) The standard error of the gap is 2 percentage
points. What do you make of the portfolio in (d)?
\end{ex}

\begin{ex}\label{ex:13.2}
An asset has a volatility of 20\% a year. You have 25 years of monthly data (300 months). Compute the standard error
of (a) its expected return, (b) its volatility, (c) its correlation with another asset, if that correlation is 0.5.
(d) You switch to daily data. Which of the three answers change?
\end{ex}

\begin{ex}\label{ex:13.3}
You optimize over 10 assets using 20 years of data. (a) If the true best Sharpe ratio is 1.0, what Sharpe ratio do you
expect to see in sample, and what do you expect to earn? (b) Redo (a) if the true best Sharpe ratio is 0.5. (c) In which
of the two worlds would you rather use $1/N$, and what do you need to know about $1/N$ to decide?
\end{ex}

\begin{ex}\label{ex:13.4}
Three assets have volatilities of 10\%, 20\% and 40\%. They are uncorrelated. Compute the weights, scaled to add to one,
under (a) risk parity, (b) minimum variance, (c) $1/N$. (d) For each rule, say what it assumes about the three expected
returns. (e) The three assets are a government bond fund, a stock index and an emerging-market fund. Which assumption
would you be least willing to make?
\end{ex}

\begin{ex}[concept]\label{ex:13.5}
For each, \tfd
(a) ``Over 1980--2026 the mean-variance portfolio of the six factors has a higher Sharpe ratio than $1/N$, so it is the
better rule.''
(b) ``$1/N$ beats mean-variance out of sample.''
(c) ``The main reason optimized weights are unstable is that volatilities are hard to estimate.''
(d) ``A long-only constraint can only hurt, since it removes portfolios the optimizer wanted.''
\end{ex}

\begin{ex}[code]\label{ex:13.6}
\texttt{F} has one row per month of factor returns. An AI wrote this loop to test mean-variance out of sample:
\begin{verbatim}
for t in range(120, len(F)):
    past = F.iloc[:t+1]
    w = np.linalg.inv(past.cov()) @ past.mean()
    oos.append(F.iloc[t] @ w)
\end{verbatim}
(a) Which months go into the estimate that earns month \texttt{t}? (b) Fix it. (c) It is one month out of hundreds. Why
does it matter so much? (d) In class the correct version had a Sharpe ratio of 0.83 and $1/N$ had 0.92. What did the
wrong version give, and what would you have concluded?
\end{ex}

\begin{ex}[concept]\label{ex:13.7}
You try five shrinkage intensities on the out-of-sample race and report the best one, which beats plain mean-variance
by 0.3 of Sharpe ratio. (a) What is wrong with the 0.3? (b) Describe how to choose the intensity without the future.
(c) Which lecture's lesson is this?
\end{ex}

% =====================================================================================================
\section{Conditional strategies I: timing risk}
{\small\emph{From: Conditional Strategies I}}

\head{The ideas}
\begin{itemize}
\item The formula $x_{t} = \mu_{t}/(\gamma\sigma_{t}^{2})$ has two parts that can move. Forecasting $\sigma$ is easy;
forecasting $\mu$ is hard. Here $\sigma_{t}$ moves and $\mu$ is held fixed.
\item Volatility comes in waves. The market's volatility is about 11\% a year in a typical month and was 93\% in March
2020. Last month's volatility explains half of this month's (in logs). It explains none of this month's return.
\item Realized variance, from daily returns within the month, is a crude forecast of next month's variance and hard to
beat by much.
\item Holding less after a volatile month pays only if the expected return does not rise with the risk. Sort the months by
last month's volatility: for the market, volatility goes from 11\% to 30\% and the average return only from 7\% to 11\%, so
the Sharpe ratio falls from 0.62 to 0.37.
\item It is a tool, not a law. It pays a lot for momentum (0.51 to 0.91) and for the best mix of the six factors, some
for profitability, nothing for the market since 1963, and it hurts for CMA, whose return does rise with its risk.
\item The worst case is a fast round trip: a crash out of calm (October 1987, still half invested) or a rebound after a
storm (April 2020, out of the market, missed $+13.6\%$).
\item The price is leverage and trading. The weight is above 2 in 16\% of months. Capping it costs little, and scaling by
$1/\sigma$ instead of $1/\sigma^{2}$ does about as well.
\end{itemize}

\begin{keybox}
\[ x_{t} = \frac{1}{\gamma}\frac{\mu_{t}}{\sigma_{t}^{2}} \qquad
   \RV_{t} = \sum_{d=1}^{N_{t}}(r_{d} - \bar r_{t})^{2} \qquad \hat\sigma_{t} = \sqrt{12\times\RV_{t}} \ \text{ a year} \]
\[ \text{if } \mu_{t} = a + b\,\sigma_{t}^{2}:\qquad x_{t} = \frac{1}{\gamma}\left(\frac{a}{\sigma_{t}^{2}} + b\right)
   \qquad b = 0:\ x_{t} \propto 1/\sigma_{t}^{2} \qquad a = 0:\ x_{t} \text{ constant} \]
\[ \text{the managed strategy:}\quad r^{\sigma}_{t} = \frac{c}{\RV_{t-1}}\,r_{t}, \qquad c \text{ such that }
   \operatorname{sd}(r^{\sigma}) = \operatorname{sd}(r) \]
\[ \text{the test:}\quad r^{\sigma}_{t} = \alpha + \beta r_{t} + \varepsilon_{t} \qquad \text{(does managing add alpha to the factor itself?)} \]
\end{keybox}

\head{Worked example}
A month has 21 trading days. If the daily returns have a standard deviation of 1\%, then
$\RV \approx 21\times 0.01^{2} = 0.0021$, and the volatility a year is $\sqrt{12\times 0.0021} = 15.9\%$. If the daily standard
deviation is 3\%, $\RV$ is nine times larger and the volatility is 47.6\%. The managed weight after the second month is
one ninth of the weight after the first.

Does that make sense? After the calmest months the market earned about 7\% with a volatility of 11\%:
$\mu/\sigma^{2} = 0.07/0.0121 = 5.8$. After the most volatile, 11\% with 30\%: $0.11/0.09 = 1.2$. The formula says to
hold almost five times as much after a calm month. Not nine times, because the return did rise a little.

\head{Exercises}

\begin{ex}\label{ex:14.1}
Suppose $\mu_{t} = a + b\,\sigma_{t}^{2}$ with $a = 4\%$ and $b = 2$, and $\gamma = 2$. (a) Compute the expected return and
the weight $x_{t}$ when volatility is 10\% and when it is 30\%. (b) Redo the two weights with $b = 0$ and $a = 4\%$.
(c) Redo them with $a = 0$ and $b = 2$. (d) In which of the three worlds does volatility timing help the most? In which
is it pointless?
\end{ex}

\begin{ex}\label{ex:14.2}
A factor is sorted into months that follow calm months and months that follow volatile months.
\begin{center}
\begin{tabular}{lcc}
\toprule
 & after calm months & after volatile months \\
\midrule
average return a year & 10\% & 12\% \\
volatility a year & 8\% & 24\% \\
\bottomrule
\end{tabular}
\end{center}
(a) Compute the Sharpe ratio in each group. (b) Compute $\mu/\sigma^{2}$ in each group. How many times larger should
the position be after a calm month? (c) Is this factor closer to $b = 0$ or to $a = 0$? (d) A second factor has a
Sharpe ratio of 0.3 after calm months and 0.6 after volatile ones. What does volatility management do to it?
\end{ex}

\begin{ex}\label{ex:14.3}
In one month the market's 21 daily returns have a standard deviation of 2\%. (a) Compute that month's realized variance
(approximately) and the volatility a year. (b) The month before, the daily standard deviation was 1\%. By what factor
does the managed weight change? (c) The managed strategy's weight averages 1. Its weight was 0.01 in April 2020, when
the market rose 13.6\%. What did the managed strategy earn that month, and what general weakness does it show?
\end{ex}

\begin{ex}[concept]\label{ex:14.4}
For each, \tfd
(a) ``Volatility management works because high volatility predicts low returns.''
(b) ``The constant $c$ is computed from the whole sample, so the comparison of Sharpe ratios looks ahead.''
(c) ``Since the volatility-managed market had an alpha with $t = 3.1$ over 1926--2026, it will raise the Sharpe ratio of
any factor.''
(d) ``Rebalancing the managed position every day on the last five days of data should do better than monthly, because
the forecast is fresher.''
\end{ex}

\begin{ex}[code]\label{ex:14.5}
\texttt{M} has one row per month with the market's return \texttt{r} and its realized variance \texttt{rv}. An AI wrote:
\begin{verbatim}
1  w = 1 / M['rv']
2  c = M['r'].std() / (w * M['r']).std()
3  managed = c * w * M['r']
\end{verbatim}
(a) One line makes this a strategy nobody could have run. Which one, and how do you fix it? (b) In which direction do
you expect the mistake to move the Sharpe ratio? Explain using the fact that the most volatile months tend to be the
months of the largest losses. (c) Line 2 uses the whole sample. Is that a problem?
\end{ex}

\begin{ex}[concept]\label{ex:14.6}
(a) Name the two costs of running the managed strategy that the Sharpe ratio does not show. (b) Give two changes that
reduce them and say what they cost in the data. (c) Which practices from \emph{Capital Allocation I} are versions of
the same idea?
\end{ex}

% =====================================================================================================
\section{Conditional strategies II: timing expected returns}
{\small\emph{From: Conditional Strategies II}}

\head{The ideas}
\begin{itemize}
\item Now the other half of the formula: let $\mu_{t}$ move. The natural signals are prices against fundamentals. From
Gordon, $D/P = r - g$: a low price against dividends means high returns ahead or low growth.
\item The dividend yield looks strong at long horizons. Over 1927--2024 the usual $t$-statistic goes from 1.8 for next
month to 12.3 for the next five years, with an $R^{2}$ of 12\%. The reason is overlap: two five-year returns a month
apart share 59 of their 60 months. A century is about 20 independent five-year returns. Allow for the overlap and
the $t$ is 2.8; use only periods that do not overlap and it is 1.5.
\item A small $R^{2}$ is not the problem. If it were real, a monthly $R^{2}$ of 1\% would raise a Sharpe ratio of 0.40 to
0.53. Whether it is real is the problem.
\item The check: each month, fit on the past only, forecast next month, and compare with the \emph{historical average},
not with zero. For D/P the out-of-sample $R^{2}$ is negative ($-0.5\%$), and timing on it earned a Sharpe ratio of 0.46
against 0.52 for buy and hold. It was short the market in 36\% of months.
\item Why it failed: the signal's average moved. Buybacks replaced dividends and D/P fell from 4\% to 1\%; the regression
read every year of that as \emph{sell}, through two of the best decades for stocks. A price ratio is a good signal only
if what is in the denominator means the same thing over time.
\item Run the same test on the other famous signals (E/P, CAPE, book-to-market, the short rate, the term spread, the
default spread, VIX, realized variance, net issuance): in sample 2 of 10 clear $|t| > 1.64$; out of sample 2 beat the
historical average, by less than 0.2\% of $R^{2}$ (Goyal and Welch).
\item For factors the same idea is the value spread: how cheap is what the factor buys against what it sells. It does
a little better out of sample, and it depends on which price is in the ratio.
\item Search finds things. Try 24 signals and about one will have $|t| > 2$ by luck. In class 5 of 24 did in the first
half of the sample, and 2 of those 5 still worked in the second half. What a real signal looks like: one signal,
several factors, both halves (factor momentum).
\item Averaging the forecasts of many weak signals did better than any one of them: the errors of the single forecasts
partly cancel, and the average stays close to the historical average.
\end{itemize}

\begin{center}
\begin{tabular}{lll}
\toprule
 & timing risk & timing returns \\
\midrule
what you forecast & next month's volatility & next month's return \\
how well, $R^{2}$ & 50\% & 0.3\%, in sample \\
what must be true & the return does not rise with risk & the signal's average stays put \\
how often it broke & by factor & by decade \\
\bottomrule
\end{tabular}
\end{center}

\begin{keybox}
\[ r_{t\to t+h} = a + b\,\frac{D}{P}\Big|_{t} + u_{t+h} \qquad
   \text{independent observations} \approx \frac{\text{years of data}}{\text{horizon in years}} \]
\[ \text{Campbell and Thompson, monthly units:}\qquad \SR^{2}_{\text{timed}} \approx \frac{\SR^{2} + R^{2}}{1 - R^{2}} \]
\[ R^{2}_{\text{OOS}} = 1 - \frac{\sum(r - \text{forecast})^{2}}{\sum(r - \text{historical average})^{2}}
   \qquad\qquad x_{t} = \frac{\text{forecast}_{t}}{\gamma\,\sigma^{2}} \]
\[ \text{value spread:}\quad VS_{t} = \log\frac{B}{M}\Big|_{\text{long},\,t} - \log\frac{B}{M}\Big|_{\text{short},\,t}
   \qquad \text{lucky signals} \approx 5\%\times\text{signals tried} \]
\end{keybox}

\head{Worked example}
A monthly $R^{2}$ of 1\%, for a market with a Sharpe ratio of 0.40 a year. In monthly units the Sharpe ratio is
$0.40/\sqrt{12} = 0.115$, so $\SR^{2} = 0.0133$. Then $(0.0133 + 0.01)/0.99 = 0.0236$ a month, and a year
$\sqrt{12\times 0.0236} = 0.53$.

The overlap. With 98 years of monthly data there are about 1{,}110 five-year returns, and $98/5 \approx 19$ that do not
overlap. A $t$-statistic that treats the 1{,}110 as separate observations is far too large.

\head{What to check before you trade a timing signal}
\begin{enumerate}[itemsep=1pt]
\item \textbf{The signal is known at $t$.} Last year's dividends, last December's book value, last month's volatility.
\item \textbf{Out of sample, month by month.} Fit on the past only. Compare with the historical average.
\item \textbf{Count the overlap.} For long-horizon returns, count the independent periods.
\item \textbf{Count your tries, and prefer a reason.}
\item \textbf{Same risk, then compare.} Against buy and hold at the same average risk, after costs.
\end{enumerate}

\head{Exercises}

\begin{ex}\label{ex:15.1}
A regression of next year's market excess return on the dividend yield gives $a = -1\%$ and $b = 2.4$. The market's
volatility is 16\% and $\gamma = 3$. (a) What return do you forecast when D/P is 4\%? When it is 1.6\%? (b) What weight do
you hold in the market in each case? (c) Over twenty years D/P drifts from 4\% to 1.6\% because firms switch from
dividends to buybacks. What does your strategy do over those years, and why is that a mistake?
\end{ex}

\begin{ex}\label{ex:15.2}
You have 60 years of monthly data and regress five-year returns on a signal. (a) How many five-year returns can you
form? (b) How many of them do not overlap? (c) The regression reports a $t$-statistic of 9. Roughly, how much evidence
do you really have? (d) Name two honest ways to report the result.
\end{ex}

\begin{ex}\label{ex:15.3}
The market has a Sharpe ratio of 0.50 a year. Use the Campbell--Thompson formula to compute the Sharpe ratio of a
timing strategy if the monthly $R^{2}$ of your signal is (a) 0.5\%, (b) 2\%. (c) Why is this an upper bound on what you
would earn?
\end{ex}

\begin{ex}\label{ex:15.4}
Four months of returns, a signal's forecast for each, and the historical average known at the time:
\begin{center}
\begin{tabular}{lcccc}
\toprule
month & 1 & 2 & 3 & 4 \\
\midrule
return & 2\% & $-1\%$ & 3\% & 0\% \\
forecast & 1\% & 1\% & 2\% & $-1\%$ \\
historical average & 1\% & 1\% & 1\% & 1\% \\
\bottomrule
\end{tabular}
\end{center}
(a) Compute the out-of-sample $R^{2}$. (b) Did the signal beat the historical average? (c) How can an out-of-sample
$R^{2}$ be negative, when an in-sample $R^{2}$ never is?
\end{ex}

\begin{ex}\label{ex:15.5}
You test 40 timing signals on the first half of your sample. Seven have $|t| > 2$. (a) How many would you expect by luck
alone? (b) You follow the seven into the second half and three still have $|t| > 2$. What do you conclude about the
other four? (c) Which is more convincing: a signal that works for one factor with $t = 3.5$ in the first half, or one
that works for four factors with $t$ near 2 in both halves? Why?
\end{ex}

\begin{ex}[concept]\label{ex:15.6}
For each, \tfd
(a) ``A signal that explains 0.5\% of next month's return is too weak to be worth trading.''
(b) ``The dividend yield has a $t$-statistic of 12 for five-year returns, so the evidence is stronger at long horizons.''
(c) ``To test a forecast out of sample, check that its out-of-sample $R^{2}$ is above zero, where zero means forecasting
a return of zero every month.''
(d) ``Book-to-market means the same thing for a software firm today as for a steel firm in 1960.''
\end{ex}

\begin{ex}[concept]\label{ex:15.7}
Volatility timing and return timing use the same formula. (a) Which input does each one forecast? (b) Why is one so
much easier than the other? Use what you know about standard errors. (c) What assumption does each need, and how did
each fail in the data?
\end{ex}

\begin{ex}[code]\label{ex:15.8}
You want to test, out of sample, whether the term spread forecasts next month's market return over 1950--2024. Give the
AI instructions. Say what goes into each month's fit, what the forecast is compared with, what number it reports, and
what would make you say the signal works. Then name one thing you must check in the code it returns.
\end{ex}

\begin{ex}[concept]\label{ex:15.9}
Nine signals each forecast next month's market return. Out of sample, most do worse than the historical average. The
average of the nine forecasts does better than every single one, and better than the historical average. (a) Explain
why, using the idea of diversification. (b) The average forecast barely moves from month to month. Why is that part of
the reason it works? (c) Connect this to shrinkage in \emph{Capital Allocation III}.
\end{ex}

% =====================================================================================================
\newpage
\section*{Formula sheet}\label{sec:sheet}
\addcontentsline{toc}{section}{Formula sheet}
{\small
\head{Returns and portfolios}
\[ r_{p,t} = \sum_{i} w_{i,t-1}r_{i,t} \quad\ w^{VW}_{i} = \frac{ME_{i,t-1}}{\sum_{j}ME_{j,t-1}} \quad\
   1 + r_{1\to T} = \prod_{t}(1+r_{t}) \quad\ \text{yearly: mean}\times 12,\ \text{vol}\times\sqrt{12} \]
\[ \text{equal weights:}\quad \sigma_{p}^{2} = \frac{\sigma^{2}}{N}\big(1 + (N-1)\bar\rho\big) \to \sigma^{2}\bar\rho
   \qquad \text{Gordon:}\quad \frac{D}{P} = r - g \]

\head{One factor}
\[ \re = \alpha + \beta\,\re_{m} + \varepsilon \qquad E[\re] = \alpha + \beta E[\re_{m}] \qquad
   \sigma^{2} = \beta^{2}\sigma_{m}^{2} + \sigma_{\varepsilon}^{2} \qquad R^{2} = \frac{\beta^{2}\sigma_{m}^{2}}{\sigma^{2}} \]
\[ \SR = \frac{E[\re]}{\sigma} \qquad \IR = \frac{E[r_{p} - r_{b}]}{\sigma(r_{p} - r_{b})} \qquad \AR = \frac{\alpha}{\sigma_{\varepsilon}}
   \qquad \SR_{\max}^{2} = \SR_{m}^{2} + \AR^{2} \qquad \text{hedged: } \re - \beta\re_{m} = \alpha + \varepsilon \]

\head{Many factors, and the cross-section}
\[ E[\re] = \alpha + \sum_{k}\beta_{k}E[f_{k}] \quad\
   \text{slope: } Z = \sum_{i}w_{i}r_{i},\ \ w_{i} = \frac{x_{i} - \bar x}{\sum_{j}(x_{j} - \bar x)^{2}},\ \
   \sum w_{i} = 0,\ \ \sum w_{i}x_{i} = 1 \]
\[ \text{characteristic-adjusted: } r_{p} - X_{p}'Z \qquad
   b = \sum_{i}w_{i}\beta_{i} \qquad \sigma_{p}^{2} = b'\Omega_{f}\,b + \sum_{i}w_{i}^{2}\sigma_{\varepsilon,i}^{2} \]

\head{What data can teach you}
\[ SE(\text{mean}) = \frac{\sigma}{\sqrt{T}} \quad SE(\alpha) = \frac{\sigma_{\varepsilon}}{\sqrt{T}} \quad
   SE(\SR) \approx \frac{1}{\sqrt{T}} \quad SE(\beta) = \frac{\sigma_{\varepsilon}/\sigma_{f}}{\sqrt{N}} \quad
   SE(\sigma) \approx \frac{\sigma}{\sqrt{2N}} \quad SE(\rho) \approx \frac{1 - \rho^{2}}{\sqrt{N}} \]
\[ t = \text{ratio}\times\sqrt{T} \qquad T = \left(\frac{z}{\text{ratio} - \text{hurdle}}\right)^{2}
   \qquad z:\ 1.28\ (90\%),\ 1.64\ (95\%),\ 1.96\ (\text{two-sided}),\ 3\ (20 \text{ tries}) \]
{\centering $T$: years. $N$: observations.\par}

\head{Hiring and firing}
\[ \text{in: } \widehat{\AR}_{T} \ge a^{*} + \frac{z_{p}}{\sqrt{T}} \qquad \text{out: } \widehat{\AR}_{T} \le a^{*} - \frac{z_{p}}{\sqrt{T}}
   \qquad \frac{\pi_{T}}{1 - \pi_{T}} = \frac{\pi}{1 - \pi}\exp\!\big(T a(\widehat{\AR}_{T} - a/2)\big) \]
\[ \text{a good hire: } \sqrt{S^{2} + a^{2}} \qquad \text{an unskilled one: } \frac{S^{2}}{\sqrt{S^{2} + a^{2}}}
   \qquad \text{leverage } L:\ \text{capital}_{t} = \text{capital}_{t-1}(1 + L r_{t}) \]

\head{Capital allocation}
\[ x^{\star} = \frac{\mu}{\gamma\sigma^{2}} \qquad x^{\star}\sigma = \frac{\SR}{\gamma} \qquad W^{\star} = \frac{1}{\gamma}\Sigma^{-1}\mu
   \qquad \text{minimum variance} \propto \Sigma^{-1}\mathbf{1} \qquad \SR_{T} = \sqrt{\mu'\Sigma^{-1}\mu} \]
\[ \SR^{2}_{\text{with}} = \SR^{2}_{\text{without}} + \Big(\frac{\alpha}{\sigma_{\varepsilon}}\Big)^{2} \qquad
   w = (1+h)\frac{1}{\gamma}\Sigma^{-1}\mu - h\,b \qquad \text{price of a mandate} = \Delta\SR\times\sigma \]
\[ w_{i} \propto \frac{\alpha_{i}}{\sigma_{\varepsilon,i}^{2}} \qquad w_{i}\sigma_{\varepsilon,i} \propto \AR_{i} \qquad
   \SR^{2} = \frac{s_{1}^{2} + s_{2}^{2} - 2\rho s_{1}s_{2}}{1 - \rho^{2}} \qquad
   \SR_{\text{pool}} = \bar s\sqrt{\frac{N}{1 + (N-1)\bar\rho}} \]
\[ \IR = IC\sqrt{\text{breadth}} \qquad \alpha = (2p - 1)\times\text{move} \qquad
   \alpha_{\text{act}} = \frac{\tau^{2}}{\tau^{2} + \omega^{2}}\alpha_{\text{view}} \]
\[ w_{1} = \frac{1}{2} + \frac{\text{gap}}{4\bar\mu}\cdot\frac{1+\rho}{1-\rho} \qquad
   \SR^{2}_{\text{in sample}} \approx \SR^{2} + \frac{N}{Y} \qquad
   \frac{\SR_{\text{you earn}}}{\SR} \approx \frac{1}{\sqrt{1 + N/(Y\,\SR^{2})}} \]

\head{Timing}
\[ x_{t} = \frac{\mu_{t}}{\gamma\sigma_{t}^{2}} \quad\ \RV_{t} = \sum_{d}(r_{d} - \bar r_{t})^{2} \quad\
   \hat\sigma_{t} = \sqrt{12\,\RV_{t}} \quad\ r^{\sigma}_{t} = \frac{c}{\RV_{t-1}}r_{t} \quad\
   \mu_{t} = a + b\sigma_{t}^{2}\ \Rightarrow\ x_{t} = \frac{1}{\gamma}\Big(\frac{a}{\sigma_{t}^{2}} + b\Big) \]
\[ R^{2}_{\text{OOS}} = 1 - \frac{\sum(r - \text{forecast})^{2}}{\sum(r - \text{average})^{2}} \qquad
   \SR^{2}_{\text{timed}} \approx \frac{\SR^{2} + R^{2}}{1 - R^{2}}\ \ (\text{monthly}) \qquad
   \text{independent observations} \approx \frac{\text{years}}{\text{horizon}} \]
}\label{sec:sheetend}

% =====================================================================================================
\ifsolutions
\newpage
\section*{Solutions}
\addcontentsline{toc}{section}{Solutions}
{\small
Short answers. For the concept questions there is more than one good answer; these give the point that earns the marks.

\head{1. Returns, portfolios and weights}

\sol{ex:1.1} (a) Value-weighted: $(600\times 2 + 250\times(-4) + 100\times 10 + 50\times 20)/1000 = 2.2\%$. Equal-weighted:
$(2 - 4 + 10 + 20)/4 = 7\%$. (b) The two smallest stocks had the largest returns. Equal weights give them half the
portfolio; value weights give them 15\%. (c) The holdings are now worth 612, 240, 110 and 60, so the weights are 59.9\%,
23.5\%, 10.8\% and 5.9\%. These are the new value weights: nothing to trade. The equal-weighted investor must sell the
two winners and buy the two that lagged to get back to 25\% each.

\sol{ex:1.2} (a) $1.5\times 0.5 = 0.75$: you lost 25\%. Adding says zero. (b) Mean $12\%$, volatility $5\times\sqrt{12} =
17.3\%$. The excess return is 0.75\% a month, so the Sharpe ratio is $(0.75/5)\times\sqrt{12} = 0.52$.

\sol{ex:1.3} (a) $52/\sqrt{5000} = 0.74\%$. (b) With many stocks $\sigma_{p}^{2} \approx \sigma^{2}\bar\rho$, so
$\bar\rho = (15.5/52)^{2} = 0.09$. (c) Almost none. With 5{,}000 stocks the $\sigma^{2}/N$ part is already gone. What is left
is the part stocks share, and more stocks do not remove it.

\sol{ex:1.4} False. In our data 89\% of stocks are skewed to the right and the market is skewed to the left ($-0.83$).
Averaging washes out each stock's own large gains. What is left is what stocks do together, and what they do together
is crash.

\sol{ex:1.5} It is the only portfolio every investor can hold at the same time: add up all holdings and you get
market-value weights. And it needs no trading when prices move, and never forces a large position in a tiny stock. The
equal-weighted return ignores the cost of trading thousands of small, illiquid stocks back to equal weights every month.

\sol{ex:1.6} (a) The panel is stacked, so \texttt{shift(1)} takes the first row of each stock from the last row of the
previous stock. It needs \texttt{df.groupby('permno')['me'].shift(1)}. (b) It weights by \texttt{me} at the end of the
month in which the return is earned (and never uses \texttt{w}). Stocks that rose this month get more weight after the
fact, so the return is biased up: in class 25.5\% a year instead of 15.7\%. (c) Line 1: the first month of every stock
should be missing; check one stock by hand. Line 2: compare the \emph{mean} with a published market return. The
correlation looks fine (0.99); only the mean shows it.

\head{2. Sorts and long-short portfolios}

\sol{ex:2.1} Bottom (A, B): $-1.4\%$. Middle (C, D): $(200\times 0 + 100\times 1)/300 = 0.33\%$. Top (E, F): $2.5\%$.
(a) Yes, each group beats the one below. (b) The long-short trades stocks worth 250 out of 550. Most of the market's
value sits in the middle, where the signal says nothing.

\sol{ex:2.2} The whole spread comes from portfolio 1, the short leg, and the pattern is not monotone. Check what is
in portfolio 1: how many stocks, how large, and whether they can be sold short at a reasonable cost. Then redo the sort
with NYSE breakpoints and value weights, or on large stocks only, and report the two legs separately.

\sol{ex:2.3} Finer buckets hold more extreme values of the signal, so the gap in average returns widens. But each
bucket holds fewer stocks and is more volatile. The $t$-statistic is the mean over its standard error; past 20 buckets
the noise grows faster than the mean. That noise is also risk you would bear.

\sol{ex:2.4} False. The weights add to zero: the short leg pays for the long leg, and the T-bill rate cancels. The Sharpe
ratio is the mean of the long-short return over its standard deviation.

\sol{ex:2.5} Any five of: the signal must be last month's (lagged within each stock); rank within each month, not over
the whole sample; breakpoints from all stocks or from NYSE stocks; equal or value weights inside each decile, and if
value, last month's market value; drop missing values before ranking; which stocks are in the universe; which return
(the month after the signal); the sample period; whether ``high'' is decile 10.

\head{3. Performance evaluation with one factor}

\sol{ex:3.1} (a) $9 - 0.5\times 6 = 6\%$. (b) Systematic volatility $0.5\times 18 = 9\%$; $\sigma_{\varepsilon} =
\sqrt{15^{2} - 9^{2}} = 12\%$; $R^{2} = 81/225 = 0.36$. (c) $\SR = 9/15 = 0.60$; $\AR = 6/12 = 0.50$. (d) The market's
Sharpe ratio is $6/18 = 0.33$, so $\sqrt{0.33^{2} + 0.50^{2}} = 0.60$. (e) $t = 0.5\times\sqrt{9} = 1.5$; $(2/0.5)^{2} = 16$
years.

\sol{ex:3.2} (a) X: $11/20 = 0.55$. Y: $1/10 = 0.10$. (b) X: $\alpha = 11 - 1.2\times 8 = 1.4\%$, $\sigma_{\varepsilon} =
\sqrt{20^{2} - 19.2^{2}} = 5.6\%$, $\AR = 0.25$. Y: $\alpha = 1 + 0.375\times 8 = 4\%$, $\sigma_{\varepsilon} = \sqrt{10^{2} - 6^{2}} = 8\%$,
$\AR = 0.50$. (c) Y. The market's Sharpe ratio is 0.5, so with Y: $\sqrt{0.25 + 0.25} = 0.71$. With X: $\sqrt{0.25 + 0.0625}
= 0.56$. (d) Y's raw return is low because its negative beta costs it $0.375\times 8 = 3\%$. The investor can undo that
exposure by holding a little more of the market. What is left is 4\% of alpha on 8\% of risk, which the market cannot give
them. Most of X's return is market exposure they already have.

\sol{ex:3.3} (a) $3/0.15 = \$20$ million. (b) Sell \$0.50 of the market per dollar. $3/0.12 = \$25$ million, with \$12.5
million of the market sold. (c) Unhedged $20\times 9\% = \$1.8$ million; hedged $25\times 6\% = \$1.5$ million. (d) Dollars per
dollar of budget are $(\alpha + \beta E[\re_{m}])/\sigma$ unhedged and $\alpha/\sigma_{\varepsilon}$ hedged. Hedging lowers
expected dollars when $(\alpha + \beta E[\re_{m}])/\sigma > \alpha/\sigma_{\varepsilon}$: when the Sharpe ratio is above the
appraisal ratio. Here $0.60 > 0.50$.

\sol{ex:3.4} (a) False. $R^{2}$ is the share of the fund's variance that is market movement. A low $R^{2}$ says most of
its risk is its own. Whether to trust the alpha is a question for its $t$-statistic. (b) False, or Depends. Against the
broad market she is paid for the value tilt itself. Measure her against a value index, or against a model with a value
factor. (c) False. What a fund adds is $\alpha/\sigma_{\varepsilon}$. The same alpha with half the residual risk adds much more.

\sol{ex:3.5} (a) $1.8/\sqrt{23} = 0.38$. (b) From its own picks: about 17\% of its 19.5\% volatility was not the market,
four fifths of the variance. (c) We chose Berkshire because it is famous, so one good record out of thousands of
managers is expected by chance. And alpha against the market alone may be exposure to a factor missing from the model.

\sol{ex:3.6} (1) \texttt{r} is a total return: subtract \texttt{rf} first, or alpha absorbs the T-bill rate and is too
high. (2) No constant: without \texttt{sm.add\_constant(mkt)} there is no intercept at all, so alpha is forced to zero
(line 2 fails) and beta is biased. (3) Alpha is a mean: multiply by 12. $\sqrt{12}$ understates it by a factor of 3.5.

\head{4. Where factors come from}

\sol{ex:4.1} (a) $g = r - D/P$: 3\% for A, 7\% for B. (b) $r = D/P + g$: 8\% for A, 4\% for B. (c) Reading (b) is ``value
earns more'': same growth, and the cheap stock offers the higher return. Reading (a) is ``cheap for a reason'': same
return, and the cheap stock is the one expected to grow slowly.

\sol{ex:4.2} (a) July 2024 through May 2025: eleven months, $t-12$ to $t-2$. (b) June 2025, because last month's winners
tend to reverse. (c) The signal would contain the return you are trying to earn: look-ahead, and enormous fake profits.
(d) A negative spread: short-term reversal.

\sol{ex:4.3} (a) P: $6 - 0.75\times 8 = 0$. Q: $5 + 0.25\times 8 = 7\%$. (b) P adds nothing: three quarters of a unit of
the market earns the same 6\% with less risk. Q raises the investor's squared Sharpe ratio by $(7\%/\sigma_{\varepsilon})^{2}$.
(c) P, which is why average returns are not the test.

\sol{ex:4.4} (a) High. The inflow pushes up the prices of the stocks the strategy buys and pushes down the ones it
sells. (b) Lower, possibly zero: the prices are now closer to right. (c) If it is pay for risk it should persist, since
whoever holds it still bears the risk. It can shrink if more investors are willing to bear it.

\sol{ex:4.5} (a) A portfolio is what you own. Both signals put the same kind of stocks, small and thinly traded, in the
extreme deciles, even though the two numbers are not related one for one across all stocks. (b) 0.89. (c) An average
near zero hides blocks. Inside a family the average absolute correlation is about 0.6, and positive and negative
correlations cancel in the average. Count families, not strategies.

\head{5. Multi-factor models}

\sol{ex:5.1} (a) Market $0.9\times 8 = 7.2$; SMB $0.5\times 2 = 1.0$; HML $-0.5\times 4 = -2.0$; RMW $0.5\times 3 = 1.5$.
Total 7.7\%. (b) $10 - 7.7 = 2.3\%$. (c) $10 - 7.2 = 2.8\%$, which is $2.3 + 1.0 - 2.0 + 1.5$. (d) Per dollar: sell \$0.90
of the market, sell \$0.50 of SMB, \emph{buy} \$0.50 of HML, sell \$0.50 of RMW. You expect 2.3\%, with only the fund's
own risk left.

\sol{ex:5.2} (a) $\sum x^{2} = 20$, so $w = (0.15, 0.05, -0.05, -0.15)$. (b) $0.15(4) + 0.05(2) - 0.05(-1) - 0.15(-1) = 0.9\%$.
(c) The weights add to zero; $\sum w_{i}x_{i} = 0.45 + 0.05 + 0.05 + 0.45 = 1$. (d) 1\%, the average return: the
equal-weighted portfolio of the four stocks.

\sol{ex:5.3} (a) $0.4/(2/\sqrt{240}) = 0.4/0.129 = 3.1$. (b) Nowhere directly. The 4{,}000 stocks make each month's slope
a well-diversified portfolio return. The $t$-statistic comes from how that return varies across the 240 months.

\sol{ex:5.4} (a) Size $0.5(2) + 0.5(1) = 1.5$; value $0.5(-1) + 0.5(-2) = -1.5$. (b) $1.5\times(-0.5) + (-1.5)\times 1.0 =
-2.25\%$. (c) Relative to the average stock the portfolio earned $1 - 2 = -1\%$. Adjusted: $-1 - (-2.25) = +1.25\%$. Large
expensive stocks had a bad month, and this portfolio did better than its characteristics implied. (d) Its holdings and
their characteristics. Not its return history, and no betas.

\sol{ex:5.5} (a) False as stated. Alpha against which model? And alpha is what the factors do not explain, not the
return above the market. (b) False. A negative loading on a factor that paid raises alpha when the factor is added.
(c) False. It says the fund's returns move with HML. (d) False. $t(\alpha) \approx \AR\times\sqrt{\text{years}}$: the span is
what counts. It can improve a little if better betas shrink the residual.

\sol{ex:5.6} Time-series: one per strategy; factor returns; betas and alpha. Cross-sectional: one per month;
characteristics; what each characteristic paid (the factor return). A sort on book-to-market also sorts, partly, on
size, profitability and whatever else comes with cheap stocks; the regression slope holds those at zero.

\sol{ex:5.7} (a) The CAPM alpha is roughly $5\% + (-1)\times E[\text{HML}] + \ldots$ When value did well the negative loading
eats the alpha, so the fund looks ordinary or bad. (b) For each dollar of the fund, add a dollar of the value factor
(long cheap stocks, short expensive ones) to cancel the $-1$, and sell half a dollar of the market if they want a beta of
one. The package's CAPM alpha then shows the 5\%.

\head{6. Decomposing a portfolio}

\sol{ex:6.1} (a) $0.5(0.8) + 0.3(1.0) + 0.2(1.5) = 1.0$. (b) Factor: $1.0^{2}\times 15^{2} = 225$. Specific:
$0.25(400) + 0.09(900) + 0.04(1600) = 100 + 81 + 64 = 245$. (c) $\sqrt{470} = 21.7\%$. (d) $225/470 = 48\%$. (e) $100/245 = 41\%$.

\sol{ex:6.2} (a) $30/\sqrt{N}$: 15\%, 6\%, 3\%. (b) The one name alone gives $0.4\times 30 = 12\%$; the other 99 add almost
nothing (12.1\% in all). (c) What matters is $\sum w_{i}^{2}$, not the number of names. This 100-stock portfolio has the
specific risk of about six equal positions.

\sol{ex:6.3} (a) $1.64/0.65 = 2.52$. (b) Print what the weights add to after the merge. (c) The 35\% of the book with no
characteristics, the smaller and stranger names. It is now a decomposition of the large holdings only.

\sol{ex:6.4} (a) Top-down: it is the only one that works without holdings. (b) Bottom-up: it reacts to a trade at once.
(c) Characteristics: no betas needed, and it describes the portfolio as it is today.

\sol{ex:6.5} (a) No: $(1.08 - 0.78)/0.22 = 1.4$ standard errors. (b) The filing shows only US-listed stocks, which is one
part of Berkshire (not the insurance operations or the private businesses). And the regression averages over five
years in which the holdings, and the holdings' own betas, changed.

\head{7. What data can teach you, and backtesting}

\sol{ex:7.1} (a) $(1.64/0.8)^{2} = 4.2$ years. (b) $(1.64/0.3)^{2} = 30$ years. (c) Only the gap between the truth and the
hurdle counts, and it fell from 0.8 to 0.3. Years go with one over the gap squared.

\sol{ex:7.2} (a) $1.2\times\sqrt{3} = 2.1$. (b) $1/\sqrt{3} = 0.58$, and $\sqrt{(1 + 1.2^{2}/2)/3} = 0.76$. (c) No. The Sharpe
ratio is a ratio: volatility is in the numerator and the denominator and cancels. Only the length of the record
matters. (d) How many funds or strategies were there to pick from, and was this one shown because of these three
years?

\sol{ex:7.3} (a) $12/\sqrt{3} = 6.9\%$. (b) $(12/16)/\sqrt{36} = 0.125$ monthly; $(12/16)/\sqrt{756} = 0.027$ daily. (c) Only
beta. Beta is about how returns move together, and more observations show more of that. Alpha is an average, and only
more years help.

\sol{ex:7.4} (a) One: $20\times 5\%$. (b) About 3.0, so 2.6 does not pass. (c) Too strict. Six signals that correlate 0.9
are close to one test, so you ran about 15 independent tests, not 20.

\sol{ex:7.5} (a) False. A random split puts test months between training months. Out of sample means out of time.
(b) False. The standard error of a Sharpe ratio is about $1/\sqrt{T}$ whatever the volatility. (c) False. The winner of
many is partly the luckiest. In class the estimate-sample winner had a Sharpe ratio of $-0.08$ in the test sample.

\sol{ex:7.6} (a) Each alpha has a standard error of about 7\%, five times the typical alpha. Sort on noisy numbers and
the top group is mostly the industries with the luckiest noise, and luck does not repeat. (b) A beta's standard error
is 0.03, small against how much betas differ across industries, so the ranking is real. (c) No, it is noisier. A recent
year picks up something that lasts about a year, such as momentum. That is an edge that moves: a long window averages
it away.

\sol{ex:7.7} Split by date. Tune on the early years, say 1980--1992: compute the Sharpe ratio of each look-back and
choose one. Test on 1993--2000: run only the chosen look-back, once, and report its Sharpe ratio with its standard
error, its alpha and its worst drawdown. There is nothing to fit here, so two samples are enough. After you see the
test number you must not go back and change the look-back, the universe or the dates.

\sol{ex:7.8} (a) A low adjusted price in 1985 means the stock split later, and stocks split because they went up. The
sort used future returns. (b) Never use data on a date that you could not have had on that date. (c) Prices as they were
on the day for the signal; adjusted prices for computing returns.

\head{8. Hiring and firing}

\sol{ex:8.1} (a) Four years: $0.5 \pm 1.64/2$, so 1.32 and $-0.32$. Sixteen years: $0.5 \pm 1.64/4$, so 0.91 and 0.09.
(b) ``Keep watching.'' The record will sit between the two lines most of the time, however long you wait. (c) No. The
walk-away line after one year is $0.5 - 1.64 = -1.14$.

\sol{ex:8.2} (a) $\exp(3\times 1\times(0.2 - 0.5)) = \exp(-0.9) = 0.41$ times as likely. (b) Odds $2\times 0.41 = 0.81$, a
probability of 45\%. (c) Odds $4\times 0.41 = 1.63$, a probability of 62\%. (d) Fire in (b): the odds are below one. Keep in
(c).

\sol{ex:8.3} (a) Odds 4 become $4/3$: 57\%. (b) Odds 1 become $1/3$: 25\%. (c) No. The news divided both odds by three.
Where each ends up depends on where they started, and the first investor may know things about the manager's process
that the second does not.

\sol{ex:8.4} (a) $\sqrt{0.6^{2} + 0.8^{2}} = 1.0$ if true. $0.36/1.0 = 0.36$ if not. (b) $(1.64/0.8)^{2} = 4.2$ years. (c) Gain
$0.4\times 5 = 2.0$. Loss $0.24\times 4.2 = 1.0$. Hire if $2.0\,\pi > 1.0\,(1 - \pi)$: $\pi > 1/3$.

\sol{ex:8.5} False. The 5\% is the chance that a manager with no skill produces a record this good. It is not the chance
that this manager has no skill. Of 100 unskilled managers, five clear the bar, and those are the ones who come to see
you. The answer depends on how many were looked at and on your hit rate.

\sol{ex:8.6} (a) The stop in percent: run less risk and you hit it less often. Leverage does not change the appraisal
ratio. (b) The principal cannot hand over their $\gamma$, and a manager who shares the gains and not the losses is
rewarded for volatility. The stop limits the damage. (c) The stop ignores skill: it can close a good team after bad
luck. The appraisal-ratio rule ignores the size of the loss and everything you knew about the manager before the
returns.

\head{9. Why should this work?}

\sol{ex:9.1} (a) $0.7\times 0.8 = 0.56$: $-44\%$. (b) $L = 2$: $0.4\times 0.6 = 0.24$, $-76\%$. $L = 3$: $0.1\times 0.4 = 0.04$,
$-96\%$. $L = 4$: the first loss is $4\times 30\% = 120\%$ of capital: wiped out, and never sees the recovery. (c)
$0.56\times 1.25 = 0.70$: still 30\% below. (d) What capital is posted against the positions (the leverage); whether
returns are compounded or added; and the window, from peak to trough and not by calendar year.

\sol{ex:9.2} Selection and overfitting: it never existed, so do not trade it. Decay: it existed and is gone, so you are
late; do not trade it either, unless there is a reason it should persist. (c) Who is on the other side and why. A
constraint that is still in place is a reason to expect the premium to come back; an error that everyone has read
about is not.

\sol{ex:9.3} (a) Investors who want high returns and cannot or may not borrow, so they buy risky, lottery-like stocks.
A constraint. (b) More likely to last: pointing out a constraint does not remove it, while a mistake can be traded
away. (c) Borrow, since the quiet stocks must be levered to matter, and hold through the bad stretch. (d) That the risk
is real: at three times leverage the 1999 loss wipes you out before the recovery.

\sol{ex:9.4} (a) False. Risk is about when an asset pays badly, and for whom. Volatility that diversifies away earns
nothing. (b) False. The CAPM says a negative beta should earn \emph{less} than T-bills. A positive return on top of a
negative beta is alpha, a puzzle for the model. (c) False. A high expected return is a low price. The business may be
poor, or risky in a way investors dislike. (d) False. An honest backtest tells you what happened. It does not tell you
why someone sold to you, and if you cannot say, you may be the one who knows less.

\sol{ex:9.5} They differ in what they already own or in what binds them. One has a job or a business exposed to the
same risk, and the other does not. One can borrow and hold through a drawdown, and the other faces a stop or a
leverage limit. One trades small, and the other at a size where they become the investor who bears the risk.

\sol{ex:9.6} (a) Overfitting: it was the best of many in the estimate sample, so 1.4 was mostly luck. Drop it. (b) A risk
premium: the strategy is short the stocks that run in a bubble, and it lost in a bubble. Keep it, at a size and
leverage that survive the next one. (c) The economic story, and what would make you stop believing it, written
\emph{before} looking at the later data.

\head{10. Momentum}

\sol{ex:10.1} (a) $12 - 0.1\times 8 - 1.2\times 7 = 2.8\%$. It is mostly momentum under another name; check whether the 2.8\%
is statistically different from zero. (b) $10 - 0.2\times 7 = 8.6\%$: a different trade. (c) The CAPM alpha counts the
momentum premium as skill. An investor who can buy momentum cheaply should pay only for what is left after it.

\sol{ex:10.2} (a) The weight goes with one over the variance: $(1/0.5)^{2} = 4$ and $(1/2)^{2} = 0.25$. (b) Sixteen. (c) The
average return and the volatility both double. The Sharpe ratio does not change.

\sol{ex:10.3} (a) $-0.68\times 15 = -10.2\%$. (b) In a bear market the stocks that fell the most are high-beta stocks.
They are the losers, so momentum is short high beta and long low beta. (c) When the market keeps falling, the negative
beta helps a little. When it rebounds sharply, momentum loses a lot. A small gain on one side and a large loss on the
other is what the seller of a call option has.

\sol{ex:10.4} (a) Its beta is $-0.30$, so $\alpha = 7.4 + 0.30\times 8.3 = 9.9\%$. (b) It pays \emph{well} in the worst
market months ($+1.3\%$ a month in the worst third), and the winners have just gone up, which if anything makes them
safer. (c) It crashes: $-65\%$ in 1932 and $-53\%$ in 2009, in rebounds. An investor with leverage or a stop cannot hold
through that.

\sol{ex:10.5} (a) False. Against the market and UMD its alpha is 2.9\% ($t = 1.8$) with a loading of about one. It is
mostly the same trade. (b) False. Only a forecast of its volatility, and last month's realized variance will do.
(c) False for the Sharpe ratio: $c$ scales every month by the same number. It does set the leverage, which you could not
have known in advance; in real time you would target a volatility.

\sol{ex:10.6} (a) Yes. The signal ends in month $t-1$ and the return is earned in month $t$. (b) The standard signal is
lagged twice, so it ends at $t-2$ and skips the reversal month. With one lag the average return is lower: 21\% against
28\% a year in class. (c) Look-ahead: the signal would contain the return you are trying to earn.

\head{11. Capital allocation I}

\sol{ex:11.1} (a) $\mu/\sigma^{2} = 0.06/0.0225 = 2.67$. So $x = 1.33$ for $\gamma = 2$ and $0.67$ for $\gamma = 4$. (b) 20\% and
10\%, which is $0.4/2$ and $0.4/4$. (c) $20/50 = 40\%$ in the market; volatility $0.4\times 15 = 6\%$; $\gamma = \SR/\text{vol} =
0.4/0.06 = 6.7$.

\sol{ex:11.2} (a) $\mu/\sigma^{2}$: $0.072/0.0144 = 5$ and $0.064/0.0064 = 10$. Weights $1/3$ and $2/3$. (b) Mean
$7.2/3 + 2(6.4)/3 = 6.67\%$. Variance $(1/9)(144) + (4/9)(64) = 44.4$, volatility $6.67\%$. Sharpe ratio $1.0 =
\sqrt{0.6^{2} + 0.8^{2}}$. (c) $10/6.67 = 1.5$ times the tangency portfolio: 50\% in A, 100\% in B, $-50\%$ in T-bills
(borrowed). Expected return 10\%. (d) $8/6.67 = 1.2$ times: volatility 8\%. (e) Both 1.0. Scaling multiplies the mean and
the volatility by the same number. (f) $W = (5, 10)/5 = (1, 2)$: three times the tangency portfolio, volatility 20\%,
which is $\SR/\gamma = 1/5$.

\sol{ex:11.3} (a) $\sqrt{0.9^{2} + 0.3^{2}} = 0.95$. (b) 0.9. It is spanned: what you hold already reproduces its return.
(c) $0.2 \pm 1.2$ is as consistent with $+2.6\%$ as with zero. And it is one sample, in sample, before costs. ``Adds
nothing here'' is not ``drop it forever''.

\sol{ex:11.4} (a) $(1 + 3)\times 0.30 = 120\%$: all of the account and 20\% borrowed. (b) $1.20 - 3\times 0.20 = 60\%$.
(c) $h = 0$: 30\%. (d) You cannot borrow against your future wages. It binds most for the young, whose human capital
is many times their savings.

\sol{ex:11.5} (a) $(0.80 - 0.65)\times 12\% = 1.8\%$ a year. (b) 1.8 times the best portfolio, with 80\% of their wealth
borrowed. Expected return $0.80\times 18\% = 14.4\%$. (c) Hold riskier assets instead of more of the best mix: move along
the curve, away from the tangency portfolio. The Sharpe ratio falls, so they earn less than 14.4\% for the same 18\%.

\sol{ex:11.6} (a) False. They hold less of the same portfolio, and more T-bills. (b) False. It never looks at $\mu$.
(c) False. Her pay already moves with technology, so her account should hold less of it. Knowing a sector is not a
reason to own more of a risk you already carry. (d) Depends. Judged with the job it comes with, an account below the
line can be exactly right. Judged alone, it looks like a mistake.

\sol{ex:11.7} (a) No. It answered the question you asked. $\gamma = 1$ says you barely mind variance: a volatility of
$\SR/\gamma = 57\%$. Nobody means that; ask what loss they could stand. (b) The principal cannot observe or hand over a
$\gamma$. A stop is observable, and it limits a manager who shares the gains and not the losses. (c) The one with the
Sharpe ratio of 2. With a higher return per unit of risk, the same volatility hits the stop less often.

\head{12. Capital allocation II}

\sol{ex:12.1} (a) 0.50, 0.50, 0.25. (b) $\alpha/\sigma_{\varepsilon}^{2}$: 12.5, 4.2, 2.1, so the first has the largest
weight. In volatility the first two are equal and the third gets half as much. (c) $\sqrt{0.25 + 0.25 + 0.0625} = 0.75$.
(d) Volatility: $10\%\times\AR_{i}/0.75$, so 6.67\%, 6.67\%, 3.33\%. Positions: divide by $\sigma_{\varepsilon}$: 1.67, 0.56
and 0.28 dollars per dollar of capital. (e) $\sqrt{0.75^{2} + 0.5^{2}} = 0.90$. Exact because the book is hedged, so it is
uncorrelated with the market.

\sol{ex:12.2} (a) $0.6\sqrt{2} = 0.85$. (b) $0.6\sqrt{2/1.5} = 0.69$. (c) $0.6\sqrt{2/0.5} = 1.2$. (d) The negatively
correlated one. All three have the same Sharpe ratio, and they are worth very different amounts.

\sol{ex:12.3} (a) $\sqrt{16}\times 0.5 = 2.0$. (b) $0.5\sqrt{16/(1 + 15\times 0.1)} = 1.26$. (c) $0.5/\sqrt{0.1} = 1.58$. (d) A
pod that is uncorrelated with the others, more than one with a high Sharpe ratio of its own.

\sol{ex:12.4} (a) $(0.55 - 0.45)\times 6\% = 0.6\%$. (b) $32\%\times\sqrt{21/252} = 9.2\%$; $IC = 0.6/9.2 = 0.065$. (c)
$0.065\times\sqrt{12} = 0.23$. (d) Volatility given: $0.23\times 10\% = 2.3\%$. Position: $2.3/32 = 7\%$. (e) You earn that
only if you have another idea like it every month.

\sol{ex:12.5} (a) $4/(4 + 16) = 0.2$, so 1\%. (b) $4/(4 + 4) = 0.5$, so 2.5\%. (c) Very unsure: the alpha you act on goes
to zero and so does the position. Certain: you act on the full 5\%.

\sol{ex:12.6} (a) $0.05\times\sqrt{400} = 1.0$. (b) $0.10\times\sqrt{4} = 0.2$. (c) The bets are not independent, so the
breadth is far below 400 and the information ratio is far below 1.

\sol{ex:12.7} (a) False. The volatility you give goes with $\alpha/\sigma_{\varepsilon}$, and the weight with
$\alpha/\sigma_{\varepsilon}^{2}$. (b) Depends on the correlation. Uncorrelated it adds a little ($\sqrt{1.5^{2} + 0.3^{2}} =
1.53$); negatively correlated it can add a lot; positively correlated enough, you would short it. (c) False. An average
near zero can hide pairs that are one trade under two names; the average absolute correlation was 0.20. (d) False. No
view is $\alpha = 0$, and the formula acts on it.

\sol{ex:12.8} (a) The Tech fund is a factor bet. The hedged pair is an alpha bet. (b) The pair by the formula, by its
appraisal ratio. The Tech fund less than the formula says, because her job already holds Tech. (c) Tech fund: that
Tech earns a premium for bearing risk. The pair: that the market is wrong about these two firms, and she knows why.
(d) Long history for the factor's premium. Recent data for the risk of both. For the alpha, the recent past and her own
judgment, since an edge decays.

\head{13. Capital allocation III}

\sol{ex:13.1} $w_{1} = 0.5 + \frac{0.01}{0.24}\times\frac{1+\rho}{1-\rho}$. (a) 54\%. (b) 62.5\%. (c) 87.5\%. (d) 129\%.
(e) Asset 2 is shorted when the ratio of the two means, $5.5/6.5 = 0.85$, is below $\rho$. So at 0.9, not at 0.8.
(f) The gap is half a standard error from zero. The portfolio is 129\% in one asset and short the other on the strength
of a number you cannot tell from zero.

\sol{ex:13.2} (a) $20/\sqrt{25} = 4\%$. (b) $20/\sqrt{600} = 0.8\%$. (c) $(1 - 0.25)/\sqrt{300} = 0.04$. (d) (b) and (c) shrink
by about $\sqrt{21}$. (a) does not change.

\sol{ex:13.3} $N/Y = 0.5$. (a) See $\sqrt{1 + 0.5} = 1.22$. Earn $1/\sqrt{1.5} = 0.82$. (b) See $\sqrt{0.25 + 0.5} = 0.87$.
Earn $0.5/\sqrt{1 + 0.5/0.25} = 0.29$. (c) In (b): the noise is twice the signal and the optimizer keeps under 60\% of a
small number. You need the Sharpe ratio $1/N$ gets from the same assets; if it is above 0.29, use $1/N$.

\sol{ex:13.4} (a) $1/\sigma$: $(4, 2, 1)/7$, so 57\%, 29\%, 14\%. (b) $1/\sigma^{2}$: $(16, 4, 1)/21$, so 76\%, 19\%, 5\%.
(c) A third each. (d) Risk parity: the same Sharpe ratio, so the mean rises with the volatility. Minimum variance: the
same mean. $1/N$: the same Sharpe ratio and the same volatility, which is false here. (e) The same mean. Bonds and
emerging markets do not earn the same return. They may earn about the same Sharpe ratio.

\sol{ex:13.5} (a) False. In sample the optimizer wins by construction. (b) Depends on the assets. It wins on similar
assets and loses badly on very different ones. (c) False. Volatilities are the best-measured input. The damage comes
from the means, and in large books from the correlations. (d) False. With estimated inputs a constraint acts as a
prior and can help out of sample.

\sol{ex:13.6} (a) Rows 0 to \texttt{t}, including month \texttt{t} itself. (b) \texttt{past = F.iloc[:t]}. (c) The weights
are tilted toward whatever did well in month \texttt{t}, every single month, and the optimizer magnifies small
differences in means. (d) 1.03, above $1/N$: the opposite conclusion.

\sol{ex:13.7} (a) The best of five was chosen after seeing the test results, so the 0.3 includes the search. It is
in-sample fitting again. (b) Each month, pick the intensity that did best on the months before, use it for the next
month, and report the track record of that procedure. (c) \emph{Backtesting}: estimate, tune, test.

\head{14. Conditional strategies I}

\sol{ex:14.1} (a) At 10\%: $\mu = 4\% + 2(0.01) = 6\%$, $x = 0.06/(2\times 0.01) = 3$. At 30\%: $\mu = 4\% + 2(0.09) = 22\%$,
$x = 0.22/(2\times 0.09) = 1.22$. (b) $0.04/0.02 = 2$ and $0.04/0.18 = 0.22$. (c) $b/\gamma = 1$ in both. (d) Most in (b),
where the return does not rise with risk at all. Pointless in (c), where it rises one for one.

\sol{ex:14.2} (a) $10/8 = 1.25$ and $12/24 = 0.50$. (b) $0.10/0.0064 = 15.6$ and $0.12/0.0576 = 2.1$: about 7.5 times.
(c) Close to $b = 0$: the variance is nine times higher and the return is only 1.2 times higher. (d) It hurts. The
strategy cuts the position exactly when the factor's Sharpe ratio is highest.

\sol{ex:14.3} (a) $\RV \approx 21\times 0.02^{2} = 0.0084$; volatility $\sqrt{12\times 0.0084} = 31.7\%$. (b) A quarter. (c)
$0.01\times 13.6\% \approx 0.1\%$. It bets that volatility persists, so a fast rebound after a storm is its worst case.

\sol{ex:14.4} (a) False. It works when high volatility does not predict \emph{higher} returns. The return can be
unchanged; it only has to rise less than the variance. (b) False. $c$ multiplies every month by the same number, so no
Sharpe ratio or $t$-statistic changes. (c) False. It helps momentum and profitability, does nothing for several
factors, and hurts CMA. For the market it depends on the period. (d) False. Five days is too little history: you react
to a few wild days that often reverse. In class the Sharpe ratio for the market fell to 0.05.

\sol{ex:14.5} (a) Line 1 uses this month's realized variance, which you only know at the end of the month. Use last
month's: \texttt{1 / M['rv'].shift(1)}. (b) Up. With this month's variance you hold almost nothing in exactly the
months of the largest losses, after the fact. (c) No. It is one constant; it sets the size, not the Sharpe ratio.

\sol{ex:14.6} (a) Leverage, with borrowing at more than the T-bill rate, and trading: the weight moves by about three
quarters of the position a month. (b) Cap the weight (at 1 or 2 it cost nothing in the data), and scale by $1/\sigma$
instead of $1/\sigma^{2}$ (0.56 against 0.52). (c) Volatility targets, risk budgets and the stop-loss.

\head{15. Conditional strategies II}

\sol{ex:15.1} (a) $-1 + 2.4\times 4 = 8.6\%$ and $-1 + 2.4\times 1.6 = 2.8\%$. (b) $0.086/(3\times 0.0256) = 1.12$ and
$0.028/0.0768 = 0.37$. (c) It sells the market year after year, from 112\% to 37\%. The yield fell because of how firms
pay out, not because expected returns fell. The signal's average moved, and the regression assumes it does not.

\sol{ex:15.2} (a) $720 - 59 = 661$. (b) $60/5 = 12$. (c) About twelve observations' worth, not 661. A $t$ of 9 is far too
large. (d) Standard errors that allow for the overlap (HAC), or the regression on the twelve periods that do not
overlap.

\sol{ex:15.3} Monthly $\SR^{2} = 0.25/12 = 0.0208$. (a) $(0.0208 + 0.005)/0.995 = 0.0260$; a year $\sqrt{12\times 0.0260} =
0.56$. (b) $(0.0208 + 0.02)/0.98 = 0.0417$; a year $0.71$. (c) It assumes the forecast is right on average and ignores
the error in estimating it. Out of sample that error is the whole problem.

\sol{ex:15.4} (a) Forecast errors: $1, -2, 1, 1$, squares add to 7. Errors of the average: $1, -2, 2, -1$, squares add to
10. $R^{2}_{\text{OOS}} = 1 - 7/10 = 0.30$. (b) Yes. (c) In sample the regression is fitted to minimize the same errors
it is scored on, so it cannot do worse than the mean. Out of sample the coefficients come from the past, and a noisy
forecast can do worse than a constant.

\sol{ex:15.5} (a) $40\times 5\% = 2$. (b) They were most likely the first half's luck. (c) The second. It was found once
and confirmed on data it was not found in, across several assets. A single $t$ of 3.5 out of 40 tries does not clear
a Bonferroni bar of about 3.2 by much, and has no second sample behind it.

\sol{ex:15.6} (a) False. If it is real, 0.5\% a month takes a Sharpe ratio of 0.50 to 0.56. The question is whether it
is real. (b) False. The observations overlap; a century holds about 20 independent five-year returns. Allowing for
that, the $t$ is 2.8, and 1.5 on periods that do not overlap. (c) False. The benchmark is the historical average
return, not zero. (d) False. Software firms expense what they build (research, software, brands), so their book value
understates what they own. The denominator has not meant the same thing over time.

\sol{ex:15.7} (a) Volatility timing forecasts $\sigma_{t}$. Return timing forecasts $\mu_{t}$. (b) The standard error of a
volatility shrinks with the number of observations, so a month of daily data is enough. The standard error of a mean
shrinks only with years. (c) Volatility timing needs the return not to rise with risk; it failed for CMA, and in some
periods for the market. Return timing needs the signal's average to stay put; D/P's did not.

\sol{ex:15.8} For each month from 1950: fit the regression of next month's return on the term spread using only pairs
that were complete before that month; forecast the month; store the forecast, the return, and the average return up
to then. Report $R^{2}_{\text{OOS}}$ against the historical average, by sub-period too, and the Sharpe ratio of a
strategy that holds forecast$/(\gamma\sigma^{2})$ against buy and hold at the same risk. It works if
$R^{2}_{\text{OOS}}$ is positive in most sub-periods. Check that the month being forecast is not inside the data used
for its own fit.

\sol{ex:15.9} (a) Each forecast is the truth plus its own estimation error. The errors are not perfectly correlated, so
averaging cancels part of them, as a portfolio does with specific risk. (b) Each regression overreacts to its own
signal. The average of nine moves less, so it stays near the historical average, which was hard to beat in the first
place. (c) It is shrinkage toward a constant: you keep a little of what the data say and lean on the prior for the
rest.
}
\fi

\end{document}
